% concatenated from repaired_paper_final.tex
\documentclass{amsart}
\usepackage{graphicx} % Required for inserting images


\usepackage{amsfonts,amsthm,amsmath,amssymb,amscd,mathrsfs}
\usepackage{graphics}
\usepackage{indentfirst}
\usepackage{cite}
\usepackage{latexsym}
\usepackage[dvips]{epsfig}
\usepackage{color, xcolor}
\usepackage{bookmark}

\usepackage{lmodern}
\usepackage{enumerate}
%\usepackage{showkeys}

\usepackage{tikz}
\usetikzlibrary{arrows.meta}


%% \usepackage{leqno, showkeys}
%%%%%%%%%%%%%%%%%%%%%%%%%%%%%%%%
\setlength{\paperheight}{11in}
\setlength{\paperwidth}{8.5in}
\addtolength{\voffset}{-0.25in}
\addtolength{\hoffset}{-0.75in}
%\setlength{\topmargin}{1in}
\setlength{\textwidth}{6.5in}
\setlength{\textheight}{8.35in}
\setlength{\footskip}{36pt}
\setlength{\marginparsep}{0pt}
\setlength{\marginparwidth}{0in}
\setlength{\headheight}{8pt}
\setlength{\headsep}{20pt}
\setlength{\oddsidemargin}{0.75in}
\setlength{\evensidemargin}{0.75in}

\renewcommand{\baselinestretch}{1.20}%\normalsize

%%%%%%%%%%%%%%%%%%%%%%%%%%%%%%%%%%%%%%%%%%%%%%
\newtheorem{theorem}{Theorem}[section]
\newtheorem{corollary}{Corollary}[section]
\newtheorem{remark}{Remark}[section]
\newtheorem{definition}{Definition}[section]
\newtheorem{lemma}[theorem]{Lemma}
\newtheorem{pro}[theorem]{Proposition}
\newtheorem{example}{Example}[section]

\newenvironment{pf}{{\noindent \it \bf Proof:}}{{\hfill$\Box$} }


%%%%%%%%%%%%%%%%%%%%%%%%%%%%%%%%%%%%%%%%%%%%%%%%%%%%
\newcommand{\red}{\color{red}}


\newcommand{\blue}{\color{blue}}

%%%%%%%%%%%%%%%%%%%%%%%%%%%%%%%%%%%%%%%%%
%%% Local Macros
%%%
\newcommand{\ti}{\tilde}
\newcommand{\lm}{\lambda}
\renewcommand{\div}{{\mathrm div \thinspace }}
\newcommand{\tn}{\tilde{\rho}}
\newcommand{\tp}{{P(\tn)}}
\newcommand{\pa}{\partial}
\newcommand{\bi}{\bibitem}
\newcommand{\rs}{\rho_s}
\newcommand{\dis}{\displaystyle}
\newcommand{\ia}{\int_0^T}
\newcommand{\bt}{\begin{theorem}}
	\newcommand{\bl}{\begin{lemma}}
		\newcommand{\el}{\end{lemma}}
	\newcommand{\et}{\end{theorem}}
\newcommand{\ga}{\gamma}
\newcommand{\curl}{\mathrm{curl}}
\newcommand{\te}{\theta}
\newcommand{\Te}{\Theta}
\newcommand{\al}{\alpha}
\newcommand{\de}{\delta}
\newcommand{\ve}{\varepsilon}
\newcommand{\la}{\label}
\newcommand{\ka}{\kappa}
\newcommand{\aaa}{\,\,\,\,\,}
\newcommand{\lz}{\,\,\,\,\,\,}
\newcommand{\ol}{\overline}
\newcommand{\bn}{\begin{eqnarray}}
	\newcommand{\en}{\end{eqnarray}}
\newcommand{\bnn}{\begin{eqnarray*}}
	\newcommand{\enn}{\end{eqnarray*}}
%\renewcommand{\theequation}{\thesection.\arabic{equation}}

%% Huang's Symbols
\newcommand{\ba}{\begin{aligned}}
	\newcommand{\ea}{\end{aligned}}
\newcommand{\be}{\begin{equation}}
	\newcommand{\ee}{\end{equation}}
\renewcommand{\O}{{\R^3}}
\newcommand{\p}{\partial}
\newcommand{\lap}{\triangle}
\newcommand{\g}{\gamma}
\renewcommand{\o}{\omega}
\newcommand{\no}{\nonumber \\}
\newcommand{\n}{\rho}
\newcommand{\si}{\sigma}
%% \newcommand{\norm[#1]#2}{\|#2\|_{#1}}
%\renewcommand{\theequation}{\thesection.\arabic{equation}}
\renewcommand{\la}{\label}
\newcommand{\na}{\nabla}
\newcommand{\on}{\bar\rho}
\newcommand{\vp}{\varphi}
%%%%%%%%%%%%%%%%%%%%%%%%%%%%%%%%%%%%%%%%%%%%%%%%%%%%%%%%%%%%%%
%%%%%%%%%%%%%%%%%%%%%%%%%%%%%%%%%%%%%%%%%%%%%%%%%%%%%%%%%%%%%%
%%%%%%%%%%%%%%%%%%%%%%%%%%%%%%%%%%%%%%%%%%%%%%%%%%%%%%%%%%%%%%
%%%%%%%%%%%%%%%%%%%%%%%%%%%%%%%%%%%%%%%%%%%%%%%%%%%%%%%%%%%%%%

\newcommand{\Bv}{{\boldsymbol{v}}}
\newcommand{\bBV}{\boldsymbol{V}}
\newcommand{\Bn}{{\boldsymbol{n}}}
\newcommand{\Bt}{{\boldsymbol{\tau}}}
\newcommand{\Bu}{{\boldsymbol{u}}}
\newcommand{\BU}{{\boldsymbol{U}}}
\newcommand{\Be}{{\boldsymbol{e}}}
\newcommand{\BF}{{\boldsymbol{F}}}
\newcommand{\Bf}{{\boldsymbol{f}}}
\newcommand{\BG}{{\boldsymbol{G}}}
\newcommand{\Bg}{{\boldsymbol{g}}}
\newcommand{\Bvphi}{{\boldsymbol{\varphi}}}
\newcommand{\supp}{\text{supp}}
\newcommand{\bBU}{\bar{{\boldsymbol{U}}}}
\newcommand{\BmcF}{{\boldsymbol{\mathcal{F}}}}
\newcommand{\BPsi}{{\boldsymbol{\Psi}}}
\newcommand{\Btau}{{\boldsymbol{\tau}}}
\newcommand{\Bnu}{{\boldsymbol{\nu}}}
\newcommand{\Bx}{{\boldsymbol{x}}}
\newcommand{\Bo}{{\boldsymbol{\omega}}}
\newcommand{\Bp}{{\boldsymbol{\phi}}}
\newcommand{\R}{\mathbb{R}}
\newcommand{\Rhalf}{\mathbb{R}^{3}_{+,*}}
\newcommand{\Sphhalf}{\mathbb{S}^{2}_{+}}
\newcommand{\msR}{\mathscr{R}}
\newcommand{\mcA}{\mathcal{A}}
\newcommand{\mcC}{\mathcal{C}}
\newcommand{\mcH}{\mathcal{H}}
\newcommand{\mcL}{\mathcal{L}}
\newcommand{\mfu}{\mathfrak{u}}
\newcommand{\mfv}{\mathfrak{v}}
\newcommand{\what}{\widehat}
\newcommand{\veps}{\varepsilon}
\newcommand{\BP}{{\boldsymbol{\Psi}}}
\newcommand{\Bb}{{\boldsymbol{b}}}
\newcommand{\Bw}{\boldsymbol{w}}
\newcommand{\ulm}{\underline{\lambda}}
\newcommand{\olm}{\overline{\lambda}}
\newcommand{\rmd}{\mathrm{d}}
%%%%%%%%%%%%%%%%%%%%%%%%%%%%%%%%%%%%%%%%%%%%%%%%%%%%%%%%%%%%%%
%%%%%%%%%%%%%%%%%%%%%%%%%%%%%%%%%%%%%%%%%%%%%%%%%%%%%%%%%%%%%%
%%%%%%%%%%%%%%%%%%%%%%%%%%%%%%%%%%%%%%%%%%%%%%%%%%%%%%%%%%%%%%%%






\title{Forced self-similar solutions to the stationary Navier--Stokes equations in a half-space}


\author{Yun Wang}
\address{School of Mathematical Sciences, Center for dynamical systems and differential equations, Soochow University, Suzhou, China}
\email{ywang3@suda.edu.cn}


\author{Chunjing Xie}
\address{School of Mathematical Sciences, Institute of Natural Sciences, Ministry of Education Key Laboratory of Scientific and Engineering Computing, and CMA-Shanghai, Shanghai Jiao Tong University, 800 Dongchuan Road, Shanghai, China}
\email{cjxie@sjtu.edu.cn}


\author{Shaoheng Zhang}
\address{School of Mathematical Sciences, Soochow University, Suzhou, China}
\email{20234007008@stu.suda.edu.cn}


\begin{document}

\begin{abstract}
We study axisymmetric self-similar solutions to the stationary Navier--Stokes equations in the half-space with the no-slip boundary condition, driven by an axisymmetric (-3)-homogeneous external force.
If the tangential curl of the force on the unit sphere is sufficiently small, we prove the existence of a unique small solution; when the force is swirl-free, the solution is automatically swirl-free and unique.
For a swirl-free external force $\boldsymbol{F}$, we introduce a scaling parameter $\lambda$ and consider the system with force $\lambda\boldsymbol{F}$; we prove that solutions exist precisely for $\lambda$ in an open interval containing zero. 
%the finiteness of the endpoints is completely determined by the sign of an explicit discriminant function.
The same approach extends to solid cones with the no-slip boundary condition, where narrower opening angles allow the existence of solutions under  larger external forces.
\end{abstract}

\maketitle



\section{Introduction}

\subsection{Problem and scaling}

We study the stationary incompressible Navier--Stokes equations
in the three-dimensional upper half-space
\[
\mathbb R^3_+
:=
\{x=(x_1,x_2,x_3)\in\mathbb R^3:x_3>0\},
\]
subject to the no-slip boundary condition:
\begin{equation}\label{eqn:NSEhfD}
	\left\{
	\begin{aligned}
		-\Delta \Bu+(\Bu\cdot\nabla)\Bu+\nabla p
		&=\BF
		&&\text{in }\mathbb R^3_+,\\
		\operatorname{div}\Bu
		&=0
		&&\text{in }\mathbb R^3_+,\\
		\Bu
		&=0
		&&\text{on }\partial\mathbb R^3_+\setminus\{0\}.
	\end{aligned}
	\right.
\end{equation}
Here $\Bu$, $p$, and $\BF$ denote the velocity, pressure, and
external force, respectively. The origin is excluded from the
boundary condition because the homogeneous solutions considered
below may have a singularity there.

The system is invariant under the scaling
\[
\Bu_\lambda(x)=\lambda\Bu(\lambda x),
\qquad
p_\lambda(x)=\lambda^2p(\lambda x),
\qquad
\BF_\lambda(x)=\lambda^3\BF(\lambda x).
\]
A solution is called self-similar if
$(\Bu_\lambda,p_\lambda)=(\Bu,p)$ for every $\lambda>0$.
Equivalently,
\[
\Bu(\lambda x)=\lambda^{-1}\Bu(x),
\qquad
p(\lambda x)=\lambda^{-2}p(x).
\]
Consequently, a scale-invariant external force must satisfy
\[
\BF(\lambda x)=\lambda^{-3}\BF(x).
\]






We study the stationary incompressible Navier--Stokes equations
in the three-dimensional upper half-space
\[
\mathbb R^3_+:=\{x=(x_1,x_2,x_3)\in\mathbb R^3:x_3>0\},
\]
subject to the no-slip boundary condition:


%Thus, a self-similar solution forces $\BF$ to be $(-3)$-homogeneous.
%In general, a function $f$ is said to be $(-k)$-homogeneous if $f(x)=\lm^k f(\lm x)$ for any $\lm>0$.

The study of self-similar solutions of the stationary Navier--Stokes equations has a long history and plays a fundamental role in understanding the local and global structures of fluid flows. Such solutions naturally appear as leading-order asymptotic descriptions of jets, wakes, and vortical motions near singularities or at large distances \cite{DG00, DI17, GPS96, JS26, KS11, MT12}. 
%The classical Landau jet~\cite{Landau44} models a round laminar jet issuing from a point source, while the swirling vortex solutions of Serrin~\cite{Serrin72} describe the interaction of a vortex filament with a plane boundary. 
More recently, self-similar solutions have provided crucial insights into the regularity theory, the Liouville-type properties, and the asymptotic behavior of solutions in various domains. In the following subsections, we briefly recall the main achievements in this area, focusing on the whole space, the half-space, and the influence of external forces.

\subsection{Related results}

Homogeneous stationary solutions play an important role in the
analysis of singularities and far-field asymptotics for the
Navier--Stokes equations. In three dimensions, the classical Landau
solutions form a one-parameter family of axisymmetric, swirl-free,
$(-1)$-homogeneous solutions in
$\mathbb R^3\setminus\{0\}$; see~\cite{Landau44}.
Their classification among axisymmetric solutions was obtained in
\cite{TianXin98,CK2004}, and \v{S}ver\'ak~\cite{Sverak11} proved
that every smooth $(-1)$-homogeneous solution in
$\mathbb R^3\setminus\{0\}$ is a Landau solution, without assuming
axisymmetry. Homogeneous axisymmetric solutions that may be singular
at the poles were further investigated in
\cite{LLY18I,LLY18II,LLY19III}.

The corresponding problem with a scale-invariant external force was
studied by Shi~\cite{Shi18}. For sufficiently small axisymmetric
$(-3)$-homogeneous forces satisfying a compatibility condition, a
family of self-similar solutions was constructed in the whole space.
In the swirl-free case, the admissible amplitude of the force is
related to the sign of an explicit discriminant. Results in higher
dimensions and for weaker notions of solutions can be found in
\cite{FR96,Tsai98,BGLWX23,BGLWX25}.

The presence of a boundary changes the problem substantially. For
the unforced stationary Navier--Stokes equations in the half-space,
Kang, Miura, and Tsai~\cite[Theorem~5.1]{KMT18} proved that every
axisymmetric $(-1)$-homogeneous solution satisfying the no-slip or
Navier-slip boundary condition is identically zero. Therefore, in
contrast with the whole-space Landau family, a nontrivial
axisymmetric homogeneous flow in the half-space must be sustained by
a non-potential external force.

\subsection{Main results for the half-space with external force}

The discussion above leads to a natural question: given a $(-3)$-homogeneous external force $\BF$, does the half-space admit a corresponding self-similar solution? 
In this paper we investigate this problem under the additional assumption that $\BF$ is axisymmetric, and we search for axisymmetric self-similar solutions.
Our first result shows that a unique axisymmetric self-similar solution exists whenever the tangential curl of the external force is sufficiently small.


\begin{theorem}\label{thm:A}
Let $\BF$ be an axisymmetric vector field of class $C^1$, homogeneous of degree $-3$ on $\mathbb{R}^3_+$, and let $(\curl\,\BF)^{\tan}$ denote the tangential component of $\curl\,\BF$ on the unit sphere $\mathbb{S}^2$.
Then there exists $\varepsilon>0$ such that if
\[
\|(\curl\,\BF)^{\tan}\|_{C^0(\mathbb{S}^2\cap\mathbb{R}^3_+)} \le \varepsilon,
\]
the system \eqref{eqn:NSEhfD} admits a unique \emph{small axisymmetric self-similar solution in the bounded angular-profile class}. More precisely,
\[
\Bu\in C^{2,\alpha}_{\mathrm{loc}}(\overline{\mathbb R^3_+}\setminus\{0\}),
\qquad p\in C^{1,\alpha}_{\mathrm{loc}}(\overline{\mathbb R^3_+}\setminus\{0\})
\]
for every $\alpha\in(0,1)$.

Moreover, if $\BF$ has no swirl, then any axisymmetric self-similar solution is necessarily swirl-free and unique.
\end{theorem}


\begin{remark}[Whole-space compatibility condition]
In the whole space $\mathbb{R}^3\setminus\{0\}$, write the axisymmetric $(-3)$-homogeneous external force in spherical coordinates as
$\BF = \rho^{-3}\bigl(X(\varphi)\Be_{\rho} + Y(\varphi)\Be_{\varphi} + Z(\varphi)\Be_{\theta}\bigr)$.
The existence of a self-similar solution 
requires the following compatibility condition:
\[
\int_0^\pi X(\varphi)\sin\varphi\cos\varphi\,\mathrm{d}\varphi
= \int_0^\pi Y(\varphi)\sin^{2}\!\varphi\,\mathrm{d}\varphi .
\]
This condition is equivalent to condition~(1.4) in \cite{Shi18}.
In the half-space (or on a solid cone), no such 
compatibility condition on the external force is needed.
\end{remark}



\begin{remark}
Theorem~\ref{thm:A} guarantees the existence of a solution when the tangential curl is small.
As a limiting case, if $\|(\curl\,\BF)^{\tan}\|_{C^0(\mathbb{S}^2\cap\mathbb{R}^3_+)}=0$, then the whole curl vanishes and the only possible self-similar solution is $\Bu\equiv0$.
Indeed, writing the axisymmetric force as $\BF=\rho^{-3}(X\Be_\rho+Y\Be_\varphi+Z\Be_\theta)$ and using \eqref{eqn:sphericalcurl},
\[
\nabla\times\BF = \frac{1}{\rho^4}\Bigl( \frac{(Z\sin\varphi)'}{\sin\varphi}\,\Be_\rho + 2Z\,\Be_\varphi - (X'+2Y)\,\Be_\theta \Bigr).
\]
Hence $\|(\curl\,\BF)^{\tan}\|_{C^0}=0$ forces $Z=0$ and $X'+2Y=0$, so $\curl\,\BF=0$.
Because $\mathbb R^3_+$ is simply connected and $\BF$ is curl-free,
there exists a $(-2)$-homogeneous scalar potential $\widetilde p$
such that $\BF=\nabla\widetilde p$(see \cite[p.~10]{Tsai18}); thus the external force can be absorbed into the pressure gradient.
The system reduces to the homogeneous Navier--Stokes equations, and by the Liouville-type result in \cite[Theorem~5.1]{KMT18} the only admissible solution under the no-slip boundary condition is $\Bu\equiv0$.
\end{remark}



The above remark shows that a vanishing tangential curl forces the whole curl to be zero.  
However, smallness of the tangential curl does \emph{not} imply smallness of the full curl: the radial component
$\frac{(Z\sin\varphi)'}{\sin\varphi}$ can be large while the tangential components remain tiny.  
%Hence the condition in Theorem~\ref{thm:A} is genuinely weaker than requiring the entire curl to be small.  
The following example exhibits a family of smooth axisymmetric forces for which the $C^{0}$-norm of the tangential curl tends to zero, yet the full curl blows up in $C^{0}$.



\begin{example}
Fix a smooth bump function $\phi\in C^{\infty}_{c}(\R)$ with $\operatorname{supp}\phi\subset[-1,1]$,
$0\le\phi\le1$, and $\phi(0)=1$.  For $\varepsilon,\delta>0$ set
\[
Z_{\varepsilon,\delta}(\varphi)=\varepsilon\,\phi\Bigl(\frac{\varphi-\frac{\pi}{4}}{\delta}\Bigr),\qquad
X_{\varepsilon,\delta}=Y_{\varepsilon,\delta}\equiv0,
\]
and take $\delta<\frac{\pi}{8}$ so that $\operatorname{supp}Z_{\varepsilon,\delta}\subset(\frac{\pi}{8},\frac{3\pi}{8})$.
Then $\BF_{\varepsilon,\delta}= \rho^{-3}Z_{\varepsilon,\delta}\,\Be_{\theta}$ is smooth and axisymmetric.
Formula \eqref{eqn:sphericalcurl} yields
\[
(\curl\,\BF_{\varepsilon,\delta})^{\tan}= \frac{2Z_{\varepsilon,\delta}}{\rho^{4}}\,\Be_{\varphi},
\qquad
\curl\,\BF_{\varepsilon,\delta}= \frac{1}{\rho^{4}}\Bigl( \frac{(Z_{\varepsilon,\delta}\sin\varphi)'}{\sin\varphi}\,\Be_{\rho}
+ 2Z_{\varepsilon,\delta}\,\Be_{\varphi} \Bigr).
\]
Since $\cot\varphi$ is bounded on the support of $Z_{\varepsilon,\delta}$,
\[
\biggl\|\frac{(Z_{\varepsilon,\delta}\sin\varphi)'}{\sin\varphi}\biggr\|_{C^{0}}
\ge \|Z_{\varepsilon,\delta}'\|_{C^{0}} - \|Z_{\varepsilon,\delta}\cot\varphi\|_{C^{0}}
\ge \frac{\varepsilon}{\delta}\|\phi'\|_{C^{0}} - C\varepsilon.
\]
Now choose $\varepsilon=\delta^{1/2}$ and let $\delta\to0^{+}$.  Then
\[
\|(\curl\,\BF_{\varepsilon,\delta})^{\tan}\|_{C^{0}(\mathbb{S}^{2}\cap\R^{3}_{+})}
= 2\|Z_{\varepsilon,\delta}\|_{C^{0}} = 2\delta^{1/2}\to0,
\]
while
\[
\|\curl\,\BF_{\varepsilon,\delta}\|_{C^{0}(\mathbb{S}^{2}\cap\R^{3}_{+})}
\gtrsim \delta^{-1/2}\to\infty.
\]
Thus the tangential part of the curl can be arbitrarily small without forcing the whole curl to be small,
and Theorem~\ref{thm:A} applies to forces whose full curl may be large.
\end{example}


Theorem~\ref{thm:A} proves the existence of a unique self-similar solution for sufficiently small forces, which raises the question of what happens when the force is large or its magnitude varies. To understand the dependence on the force magnitude in the swirl-free case, we introduce a parameter $\lambda$ and replace $\BF$ by $\lambda\BF$.  The following theorem shows that solutions exist exactly for $\lambda$ in an open interval $(\underline{\lambda},\overline{\lambda})$ containing $0$, and the finiteness of the endpoints is determined by the sign of an explicit discriminant.


\begin{theorem}
\label{thm:B}
Let $\BF$ be an axisymmetric, swirl-free external force on $\mathbb{R}^{3}_{+}$, of class $C^{1}$ and homogeneous of degree $-3$.
There exist constants $\underline{\lambda}<0<\overline{\lambda}$ (possibly $\pm\infty$) with the following property: for every $\lambda\in(\underline{\lambda},\overline{\lambda})$ the Navier--Stokes equations \eqref{eqn:NSEhfD} with external force $\lambda\BF$ admit a unique axisymmetric self-similar solution
$\Bu\in C^{2,\alpha}_{\mathrm{loc}}(\overline{\mathbb R^3_+}\setminus\{0\})$, $p\in C^{1,\alpha}_{\mathrm{loc}}(\overline{\mathbb R^3_+}\setminus\{0\})$ for every $\alpha\in(0,1)$, whereas no such bounded-profile solution exists for $\lambda\notin(\underline{\lambda},\overline{\lambda})$.

In spherical coordinates we write $\BF = \rho^{-3}\bigl( X(\varphi)\Be_{\rho} + Y(\varphi)\Be_{\varphi} \bigr)$.
Then the finiteness of $\underline{\lambda}$ and $\overline{\lambda}$ is determined by the sign of the discriminant
\begin{align}\label{eqn:disfcn}
\mathcal{N}(\vp)=\int_{0}^{\vp}
\bigl(G(\tau)+2\mathcal{C}\bigr)\,\sin\tau\,(\cos\vp-\cos\tau)\,\rmd\tau,
\quad \vp \in \big[0,\tfrac{\pi}{2}\big],
\end{align}
where $G(\vp)=X(\vp)+2\int_{0}^{\vp} Y(\tau)\,\rmd\tau$ and $\mathcal{C}=-\int_{0}^{\frac{\pi}{2}}G(\tau)\sin\tau\cos\tau\,\rmd\tau$. More precisely:
\begin{itemize}
\item if $\mathcal{N}\ge0$ and $\mathcal{N}\not\equiv0$, then $\underline{\lambda}\in(-\infty,0)$ and $\overline{\lambda}=+\infty$;
\item if $\mathcal{N}\le0$ and $\mathcal{N}\not\equiv0$, then $\underline{\lambda}=-\infty$ and $\overline{\lambda}\in(0,+\infty)$;
\item if $\mathcal{N}$ changes sign, then $\underline{\lambda}\in(-\infty,0)$ and $\overline{\lambda}\in(0,+\infty)$;
\item if $\mathcal{N}\equiv0$, then $\underline{\lambda}=-\infty$ and $\overline{\lambda}=+\infty$.
\end{itemize}
\end{theorem}


\begin{remark}[Blow-up location]
As proved in Lemma~\ref{lm:exactblow}, a finite $\underline{\lambda}$ or $\overline{\lambda}$ corresponds to a solution that blows up exactly at the north pole $\varphi=0$; that is, it cannot be extended continuously to the symmetry axis.
\end{remark}



If $\mathcal{N}\equiv0$, then \eqref{eqn:disfcn} forces $G+2\mathcal{C}\equiv0$; hence $G$ is constant, and
$\nabla\times\BF = -\rho^{-4}(X'+2Y)\Be_\theta = 0$. Therefore the force is a gradient and can be absorbed into the pressure. Then only the trivial solution exists (see \cite[Theorem 5.1]{KMT18}), corresponding to $\underline{\lambda}=-\infty$ and $\overline{\lambda}=+\infty$ in Theorem~\ref{thm:B}. For the non-trivial cases where $\mathcal{N}\not\equiv0$, the sign of $\mathcal{N}$ determines the three possible behaviours of the discriminant, which are illustrated by the following examples.


\begin{example}\label{example:N}
For $G(\varphi)=X(\varphi)+2\int_0^\varphi Y(\tau)\,\mathrm{d}\tau$ and $\mcC=-\int_0^{\frac{\pi}{2}} G(\tau)\sin\tau\cos\tau\,\mathrm{d}\tau$, the choices
\begin{itemize}
\item $G(\varphi)=\mp\cos(2\varphi)$ give $\mcC=0$ and
  $\mathcal{N}(\varphi)=\pm\frac{1}{6}\cos\varphi(\cos\varphi+2)(1-\cos\varphi)^2$, strictly positive/negative on $(0,\tfrac{\pi}{2})$;
\item $G(\varphi)=\cos(4\varphi)$ gives $\mcC=\frac16$ and
  $\mathcal{N}(\varphi)=\frac{2}{15}\cos\varphi(1-\cos\varphi)^2\bigl(-2\cos^3\varphi-4\cos^2\varphi-\cos\varphi+2\bigr)$, which changes sign.
\end{itemize}
These illustrate the three non-trivial cases in Theorem~\ref{thm:B}.
\end{example}

%\subsection{Relation to previous work}

%It is instructive to position our results within the existing literature.
%\begin{itemize}
%\item \textbf{Whole space (Shi~\cite{Shi18}).}  
 % In $\R^3\setminus\{0\}$ the existence of a self-similar solution requires the compatibility condition 
  %$\int_0^\pi X\sin\varphi\cos\varphi\,\rd\varphi = \int_0^\pi Y\sin^2\varphi\,\rd\varphi$ to fix the integration constant.
  %On the half-space the Dirichlet boundary condition at $\varphi=\pi/2$ determines this constant directly
  %(see Remark~1), thereby eliminating the compatibility restriction.
  %Moreover, while Shi's blow-up/non-existence criterion relies on the sign of a discriminant,
  %our Theorem~\ref{thm:B} gives a complete classification of the admissible amplitude interval
  %together with the exact blow-up location (\emph{cf.} Lemma~\ref{lm:exactblow}).

%\item \textbf{Half-space without force (Kang--Miura--Tsai~\cite{KMT18}).}
  %\cite[Theorem~5.1]{KMT18} proves that every $(-1)$-homogeneous axisymmetric solution
  %satisfying the Dirichlet (or Navier-slip) condition must be trivial.
 % Our results show that a $(-3)$-homogeneous external force can sustain non-trivial self-similar flows,
  %and we quantify how large $\|\BF\|$ can be before the solution ceases to exist.

%\item \textbf{Swirl vs.\ swirl-free.}
 % For swirl-free forces Shi~\cite{Shi18} proves a sharp dichotomy:
  %one direction of $\lambda$ is unbounded while the other is finite.
  %We recover the same phenomenon for the half-space (Theorem~\ref{thm:B}),
  %but our discriminant $\mathcal{N}(\varphi)$ has a transparent geometric meaning:
  %it is the tangential curl integrated against spherical-cap weight functions
  %(see Remark~\ref{rmk:discriminant-structure}).
%\end{itemize}


%Theorem~\ref{thm:A} shows that, provided the tangential curl does not exceed a universal threshold, a unique small self-similar flow exists and is controlled by the force. In the swirl-free case, Theorem~\ref{thm:B} reveals that the range of admissible force amplitudes $\lambda$ is determined by the sign of a simple discriminant $\mathcal{N}$; in particular, when $\mathcal{N}$ changes sign, the interval is bounded, indicating that too strong an external driving or a reversal of its direction destroys the self-similar solution. These mathematical findings have clear physical interpretations: a moderate non-potential force near a wall sustains an axisymmetric self-similar flow, while an excessive force or an unfavourable force distribution extinguishes the flow. Our analysis thus provides a rigorous foundation for controlling and predicting such bound flows in applications.
\subsection{Main ideas of the proofs}

For an axisymmetric self-similar velocity field, we write
\[
\Bu(\rho,\theta,\varphi)
=
\frac{1}{\rho}
\bigl(
f(\varphi)\Be_\rho
+
g(\varphi)\Be_\varphi
+
h(\varphi)\Be_\theta
\bigr).
\]
The divergence-free condition expresses the radial component
$f$ in terms of the meridional component $g$. After eliminating the
pressure, the swirl-free system can be integrated and transformed,
through
\[
t=\cos\varphi,
\qquad
L(t)=g(\varphi)\sin\varphi,
\]
into a second-order scalar equation. A further normalization reduces
it to the Riccati initial-value problem
\[
M'(t)+\frac12M(t)^2=\lambda N(t),
\qquad M(0)=0.
\]
Existence of a bounded angular profile is equivalent to continuation
of this Riccati solution up to $t=1$. Comparison arguments then show
that the admissible set of amplitudes is an open interval, while the
sign of the discriminant determines whether either endpoint is
finite.

When swirl is present, the azimuthal equation is first solved for
$h$ in terms of $g$ and the swirl component of the force. The
remaining equations define a nonlinear map for $g$. For sufficiently
small tangential curl, this map is a contraction on a bounded ball
of angular profiles. The cone case follows from the same construction
after keeping track of the dependence of the estimates on the
opening angle.


\subsection{Outline} 
The remainder of the paper is organized as follows.  
Section~\ref{sec:Pre} introduces the  spherical coordinates and basic notation.  
In Section~\ref{sec:No-Swirl} we reduce the swirl-free case to a Riccati equation and studies its maximal interval of existence.  
Section~\ref{sec:withswirl} treats forces with swirl by a contraction argument.  
Finally, Section~\ref{sec:Discussions} extends the results to solid cones and analyzes the dependence on the opening angle.



\section{Preliminaries}\label{sec:Pre}
Recall that the spherical coordinates of $\R^3$ are $\rho,\te,\vp$ with 
\[
\Bx=(\rho\sin \vp \cos \te,\rho \sin \vp \sin \te, \rho \cos \vp ),
\]
and the corresponding orthonormal basis vectors are
\begin{align} \label{eqn:sphericalbasis}
\Be_{\rho}=\frac{\Bx}{\rho}, \quad \Be_{\te}=(-\sin \te, \cos \te,0), \quad \Be_{\vp}=\Be_{\te}\times \Be_{\rho}.
\end{align}
Here $\rho=|\Bx|$, $\te \in (0,2\pi)$ is the azimuthal angle, and $\vp\in (0,\pi)$ is the polar angle.
In this coordinate system, a vector-valued function $\Bv$ admits the representation
\begin{equation}\label{eqn:sphericalrep}
\Bv = v^\rho(\rho,\vp,\te) \Be_\rho + v^\vp(\rho,\vp,\te) \Be_\vp + v^\te(\rho,\vp,\te) \Be_\te, 
\end{equation}
where $v^\theta$ is known as the {\it swirl} component. We say $\Bv$ is {\it axisymmetric} if it is of the form
\[
\Bv = v^\rho(\rho,\vp) \Be_\rho + v^\vp(\rho,\vp) \Be_\vp + v^\te(\rho,\vp) \Be_\te. 
\]

Let $\Bv$ be a vector field expressed in \eqref{eqn:sphericalrep}. Then the curl  of $\Bv$ is given by
\begin{equation}\label{eqn:sphericalcurl}
\begin{aligned}
\curl \, \Bv=
\frac{1}{\rho\sin\vp}
\left(
\frac{\p(v^{\te}\sin\vp)}{\p\vp}-\frac{\p v^{\vp}}{\p\te}
\right)
\Be_{\rho}
+\frac{1}{\rho}
\left(\frac{1}{\sin\vp}\frac{\p v^{\rho}}{\p\te}-\frac{\p(\rho v^{\te})}{\p\rho}\right)\Be_{\vp}
+\frac{1}{\rho}
\left(\frac{\p(\rho v^{\vp})}{\p\rho}-\frac{\p v^{\rho}}{\p\vp}\right)\Be_{\te},  
\end{aligned}
\end{equation}
see, for example, \cite[eq.~(2.47)]{AW95}.
If $\Bv$ is axisymmetric, then
\[
\curl\,\Bv = \frac{1}{\rho\sin\vp}\frac{\p(v^{\te}\sin\vp)}{\p\vp}\Be_{\rho}
- \frac{1}{\rho}\frac{\p(\rho v^{\te})}{\p\rho}\Be_{\vp}
+ \frac{1}{\rho}\Bigl(\frac{\p(\rho v^{\vp})}{\p\rho} - \frac{\p v^{\rho}}{\p\vp}\Bigr)\Be_{\te}.
\]
If, in addition, $\Bv$ is swirl-free, then $\curl\,\Bv$ has only the $\Be_{\theta}$ component.


\section{Analysis of solutions for a swirl-free external force}\label{sec:No-Swirl}


Assume that $\Bu$, $p$, and $\BF$ are axisymmetric and homogeneous of degrees $-1$, $-2$, and $-3$, respectively, and take the form
\begin{align*}
\Bu = \frac{f(\varphi)}{\rho} \Be_\rho + \frac{g(\varphi)}{\rho} \Be_\varphi + \frac{h(\varphi)}{\rho} \Be_\theta, \quad
p = \frac{P(\varphi)}{\rho^2},
\quad 
\BF = \frac{1}{\rho^3}\bigl( X(\varphi)\Be_\rho + Y(\varphi)\Be_\varphi + Z(\varphi)\Be_\theta \bigr),
\end{align*}
where $\Be_\rho, \Be_\varphi, \Be_\theta$ are the spherical basis vectors defined in \eqref{eqn:sphericalbasis}.  
Substituting these expressions into the Navier--Stokes equations \eqref{eqn:NSEhfD} yields
\begin{equation}\label{eqn:halfD1}
\left\{
\begin{aligned}
   f^{\prime\prime} + f^{\prime} \cot \varphi &= gf^{\prime} -(f^2 + g^2 + h^2) - 2P - X, \\
   f^{\prime} &= gg^{\prime} - h^2 \cot \varphi + P^{\prime} - Y, \\ 
   (h^\prime + h\cot \varphi)^{\prime} &= g(h^\prime + h\cot \varphi) - Z,\\
   f+ g^\prime + g\cot \varphi &=0,
\end{aligned}
\right.
\end{equation}
together with the compatibility condition (see \cite[Corollary~1, Lemma~2]{LW09})  and Dirichlet boundary condition
\begin{align}\label{eqn:halfDbc}
f^\prime(0) = g(0) = h(0) = 0, \quad f\bigl(\tfrac{\pi}{2}\bigr) = g\bigl(\tfrac{\pi}{2}\bigr) = h\bigl(\tfrac{\pi}{2}\bigr) = 0.
\end{align}
{\bf In this section we focus on the swirl-free case, that is, $Z\equiv0$.}



\subsection{Reduction for swirl-free external forces}\label{sec:reductNo-swirl}
We now show how system \eqref{eqn:halfD1} reduces to a Riccati equation in the swirl-free setting.
{\bf Claim: The solution $\Bu$ is also swirl-free, i.\,e., $h\equiv0$.} 
The proof follows the same idea as the derivation of the Landau solution (see, for example, \cite[Theorem~8.1]{Tsai18}).
Set $H = h^\prime + h\cot \varphi $. Equation \eqref{eqn:halfD1}$_3$ gives $H^\prime = gH$,
hence $H$ does not change sign on $[0,\frac{\pi}{2}]$. On the other hand, 
\begin{align*}
\int_{0}^{\frac{\pi}{2}} H \sin \varphi\, \mathrm{d}\varphi = \int_0^{\frac{\pi}{2}} (h\sin \varphi)^\prime \, \mathrm{d}\varphi = h(\tfrac{\pi}{2})= 0. 
\end{align*}
Thus $H\equiv0$, i.\,e.\ $h^\prime + h \cot \varphi =0$. With $h(0)=0$ from \eqref{eqn:halfDbc}, we obtain $h\equiv0$.
Now equations \eqref{eqn:halfD1} reduce to 
\begin{equation}\label{eqn:halfD2}
\left\{ 
\begin{aligned}
    f^{\prime \prime} + f^\prime \cot \varphi &= f^\prime g - f^2 - ( g^2+2P) -X,\\
    f^\prime &= g g^\prime + P^\prime - Y, \\
    f+ g^\prime + g\cot \varphi &=0. 
\end{aligned}
\right.
\end{equation}
From \eqref{eqn:halfD2}$_2$ we obtain
\begin{align}\label{eqn:eqnforP}
\frac12 g^2 + P = f + \int_0^{\varphi} Y(\tau)\,\mathrm{d}\tau + \mcC
\end{align}
for some constant $\mcC$.  Substituting this into \eqref{eqn:halfD2}$_1$ and using
$f''+f'\cot\varphi = (f'\sin\varphi)'/\sin\varphi$ gives
\begin{equation}\label{eqn:reduct1}
\frac{(f'\sin\varphi)'}{\sin\varphi} = f'g - f^2 - 2f - \bigl(G+2\mcC\bigr),
\end{equation}
where $G(\varphi) = X(\varphi) + 2\int_0^{\varphi} Y(\tau)\,\mathrm{d}\tau$.
With $-f\sin\varphi = (g\sin\varphi)'$ from \eqref{eqn:halfD2}$_3$, equation \eqref{eqn:reduct1} becomes
\[
-(f'\sin\varphi)' + (fg\sin\varphi)' + (2g\sin\varphi)' = \bigl(G+2\mcC\bigr)\sin\varphi .
\]
Integrating from $0$ to $\varphi$ and applying $f'(0)=g(0)=0$ from \eqref{eqn:halfDbc} yields
\begin{equation}\label{eqn:reduct2}
-f'\sin\varphi + fg\sin\varphi + 2g\sin\varphi = \int_0^{\varphi} \bigl(G(\tau)+2\mcC\bigr)\sin\tau\,\mathrm{d}\tau .
\end{equation}
Set $t = \cos\varphi$ and $L(t) = g(\varphi)\sin\varphi$. From \eqref{eqn:halfD2}$_3$ we compute
\begin{align*}
\frac{\mathrm{d} L}{\mathrm{d} t} = f,\qquad
\frac{\mathrm{d} f}{\mathrm{d} \varphi} = \frac{\mathrm{d}^2 L}{\mathrm{d} t^2}(-\sin\varphi).
\end{align*}
Throughout the paper, derivatives with respect to $\varphi$ are indicated either by a prime $'$ or by $\frac{\mathrm{d}}{\mathrm{d}\varphi}$, whereas derivatives with respect to the variable $t = \cos\varphi$ are always written as $\frac{\mathrm{d}}{\mathrm{d}t}$.
The boundary conditions \eqref{eqn:halfDbc} translate to
\[
L(0)=g\bigl(\frac{\pi}{2}\bigr)\sin\frac{\pi}{2}=0,\quad
L(1)=g(0)\sin0=0,\quad
\frac{\mathrm{d} L}{\mathrm{d} t}(0)=f\bigl(\frac{\pi}{2}\bigr)=0.
\]
Define $J(t):=\int_0^{\varphi} (G(\tau) + 2\mcC) \sin\tau \, \mathrm{d}\tau$.  Then \eqref{eqn:reduct2} becomes
\begin{align}\label{eqn:reduct3}
    (1-t^2)\frac{\mathrm{d}^2 L}{\mathrm{d} t^2} + L\frac{\mathrm{d} L}{\mathrm{d} t} + 2L = J(t),
\end{align}
The left-hand side is an exact derivative:
\[
(1-t^2)\frac{\mathrm{d}^2 L}{\mathrm{d} t^2} + L\frac{\mathrm{d} L}{\mathrm{d} t} + 2L
= \frac{\mathrm{d}}{\mathrm{d} t}\Bigl( (1-t^2)\frac{\mathrm{d} L}{\mathrm{d} t}+\frac{1}{2}L^2 +2 t L \Bigr).
\]
Integrating \eqref{eqn:reduct3} from $t$ to $1$ and using $L(1)=0$, we obtain
\begin{equation}\label{eqn:reduct4}
\left\{
\begin{aligned}
&(1-t^2)\frac{\mathrm{d} L}{\mathrm{d} t}+\frac{1}{2}L^2 +2 t L= -\int_{t}^1 J(s) \, \mathrm{d} s,\\
&L(0)=L(1)=\frac{\rmd L}{\rmd t}(0)=0.
\end{aligned}
\right.
\end{equation}
Evaluating the first equation at $t=0$ and using $L(0)=\frac{\mathrm{d} L}{\mathrm{d} t}(0)=0$ gives the compatibility condition
$\int_0^1 J(s)\,\mathrm{d}s = 0$, which explicitly is
\begin{equation}\label{eqn:constC}
\begin{aligned}
\int_0^1 J(s) \, \mathrm{d} s 
&= \int_0^{\frac{\pi}{2}} \biggl(\int_0^{\varphi} (G(\tau) +2 \mcC) \sin\tau \, \mathrm{d}\tau \biggr) \sin\varphi \, \mathrm{d}\varphi\\
&= \int_0^{\frac{\pi}{2}} G(\tau) \sin\tau \biggl(\int_{\tau}^{\frac{\pi}{2}}
\sin\varphi \, \mathrm{d}\varphi \biggr) \, \mathrm{d}\tau + \mcC \int_0^{\frac{\pi}{2}} 2\sin\tau(1-\cos\tau)\,\mathrm{d}\tau\\
&= \int_0^{\frac{\pi}{2}} G(\tau) \sin\tau \cos\tau \, \mathrm{d}\tau + \mcC = 0.
\end{aligned}
\end{equation}
Hence the constant $\mcC$ is determined by
\begin{align*}
\mcC 
&= -\int_0^{\frac{\pi}{2}} G(\tau) \sin\tau \cos\tau \, \mathrm{d}\tau
=-\int_0^{\frac{\pi}{2}} \Big[X(\varphi) \sin \varphi \cos \varphi + Y(\varphi) \cos^2 \varphi \Big] \, \mathrm{d}\varphi.
\end{align*}
Introduce $M(t) = \dfrac{L(t)}{1-t^2}$. Substituting $L = (1-t^2)M$ into the differential equation of \eqref{eqn:reduct4} we get
\[
\frac{\mathrm{d} M}{\mathrm{d} t} + \frac12 M^2 = N \quad\text{in } (0,1),
\]
where
\[
N(t):= \frac{-1}{(1-t^2)^2} \int_t^1 J(s) \, \mathrm{d} s.
\]
The boundary conditions in \eqref{eqn:reduct4} become $M(0)=\frac{\mathrm{d} M}{\mathrm{d} t}(0)=0$. However, because $\int_0^1 J=0$ implies $N(0)=0$, the ODE forces $\frac{\mathrm{d} M}{\mathrm{d} t}(0)=0$ automatically once $M(0)=0$ is imposed. Thus the problem \eqref{eqn:reduct4} reduces to the initial value problem
\begin{equation}\label{eqn:Riccati}
\left\{
\begin{aligned}
\frac{\mathrm{d} M}{\mathrm{d} t} + \frac12 M^2 &= N \quad \text{in } (0,1), \\
M(0)&=0.
\end{aligned}
\right.
\end{equation}
Note that $M(t)=\dfrac{g(\varphi)}{\sin\varphi}$ with $t=\cos\varphi$.
We verify that $N(t)$ extends continuously to $[0,1]$.  From
\[
\frac{\mathrm{d} J}{\mathrm{d} t}
= \frac{\mathrm{d} J}{\mathrm{d}\varphi}\frac{\mathrm{d}\varphi}{\mathrm{d} t}
= -\bigl(G(\varphi)+2\mcC\bigr),
\]
we obtain $\frac{\mathrm{d} J}{\mathrm{d} t}(1)=-(G(0)+2\mcC)$.  Since $J(1)=0$, expanding $-\int_t^1 J(s)\,\mathrm{d}s$ around $t=1$ gives
\[
-\int_t^1 J(s)\,\mathrm{d}s
= \frac12 \, \frac{\mathrm{d} J}{\mathrm{d} t}(1)\,(t-1)^2 + o\bigl((t-1)^2\bigr)
= -\frac{G(0)+2\mcC}{2}(t-1)^2 + o\bigl((t-1)^2\bigr).
\]
Dividing by $(1-t^2)^2=(1-t)^2(1+t)^2$ yields
\[
\lim_{t\to 1^-} N(t) = -\frac{1}{8}\bigl(G(0)+2\mcC\bigr)
= -\frac{1}{8}\bigl(X(0)+2\mcC\bigr).
\]
Thus $N\in C^0[0,1]$, and the Riccati equation \eqref{eqn:Riccati} is well-posed on the closed interval.
The above reduction is carried out under the assumption that a sufficiently regular self-similar solution exists. This will be justified in Section~\ref{sec:EUNo-swirl}, where we prove local regularity on $\overline{\mathbb R^3_+}\setminus\{0\}$; consequently all derivatives in this subsection are classical away from the origin.


\subsection{Existence and uniqueness for small swirl-free forces}\label{sec:EUNo-swirl}

We now prove Theorem~\ref{thm:A} for swirl-free forces.
To this end we first establish a general solvability result for the Riccati equation on an arbitrary interval $[t_0,1]$; this lemma will also be a key tool when we extend the analysis to solid cones in Section~\ref{sec:Discussions}.
Under the smallness condition on the external force, we show that the Riccati equation is solvable and that its solution yields a full Navier--Stokes solution with $\Bu\in C^{2,\alpha}$ and $p\in C^{1,\alpha}$ for every $\alpha\in (0,1)$.

\begin{lemma}[Solvability of the Riccati equation on a general interval]\label{lm:Riccati}
Let $t_0<1$, $\ell:=1-t_0$, and consider the initial value problem
\[
\frac{\mathrm{d}M}{\mathrm{d}t} + \frac12 M^{2} = N \quad \text{on } [t_0,1],\qquad M(t_0)=0,
\]
with $N\in C^{0}[t_0,1]$.
\begin{itemize}
\item[(i)] If $\|N\|_{C^{0}[t_0,1]} < \dfrac{1}{2\ell^{2}}$, then there exists a unique solution $M\in C^{1}[t_0,1]$ satisfying
\[
\|M\|_{C^{0}[t_0,1]} \le 2\ell\|N\|_{C^{0}[t_0,1]},\qquad
\|M\|_{C^{1}[t_0,1]} \le 2(\ell+1)\|N\|_{C^{0}[t_0,1]}.
\]
\item[(ii)] For any $N_{1},N_{2}$ with $\|N_{i}\|_{C^{0}[t_0,1]} \le \delta < \dfrac{1}{2\ell^{2}}$, let $M_{1},M_{2}$ be the corresponding unique solutions. Then
\[
\|M_{1}-M_{2}\|_{C^{1}[t_0,1]} \le \frac{\ell+1}{1-2\ell^{2}\delta}\,\|N_{1}-N_{2}\|_{C^{0}[t_0,1]}.
\]
\end{itemize}
\end{lemma}

\begin{proof}
(i) For brevity we omit the interval $[t_0,1]$ in the norm notation.
Define the operator $\Psi$ on $C^{0}[t_0,1]$ by
\[
\Psi M(t) := \int_{t_0}^{t} \Bigl( N(s) - \frac12 M^{2}(s) \Bigr) \, \mathrm{d}s ,\qquad t\in[t_0,1].
\]
For any $M\in C^{0}[t_0,1]$,
\[
\|\Psi M\|_{C^{0}} \le \ell\Bigl( \|N\|_{C^{0}} + \frac12 \|M\|_{C^{0}}^{2} \Bigr).
\]
Let $\mathfrak{X}_{R} = \{\, M\in C^{0}[t_0,1] : M(t_0)=0,\ \|M\|_{C^{0}} \le R \,\}$ with $R>0$ to be chosen.
If $M\in\mathfrak{X}_{R}$, then
\[
\|\Psi M\|_{C^{0}} \le \ell\Bigl( \|N\|_{C^{0}} + \frac12 R^{2} \Bigr).
\]
Choosing $R$ to satisfy $\ell\|N\|_{C^{0}} + \frac{\ell}{2}R^{2} = R$ gives
\[
R = \frac{1}{\ell}\bigl(1 - \sqrt{\,1 - 2\ell^{2}\|N\|_{C^{0}}\,}\bigr),
\]
which is well defined and positive because $\|N\|_{C^{0}} < \frac{1}{2\ell^{2}}$.
With this $R$, $\Psi$ maps $\mathfrak{X}_{R}$ into itself.

Now estimate the difference of two images:
\[
(\Psi M_{1} - \Psi M_{2})(t) = \frac12 \int_{t_0}^{t} \bigl( M_{2}^{2} - M_{1}^{2} \bigr) \, \mathrm{d}s,
\]
hence
\[
\|\Psi M_{1} - \Psi M_{2}\|_{C^{0}}
   \le \frac{\ell}{2} \bigl( \|M_{1}\|_{C^{0}} + \|M_{2}\|_{C^{0}} \bigr) \|M_{1}-M_{2}\|_{C^{0}}
   \le \ell R \,\|M_{1}-M_{2}\|_{C^{0}}.
\]
Because $\ell R = 1 - \sqrt{\,1-2\ell^{2}\|N\|_{C^{0}}\,} < 1$, the map $\Psi$ is a contraction on $\mathfrak{X}_{R}$. Consequently, there exists a unique $M\in\mathfrak{X}_{R}$ with $M = \Psi M$, which solves the Riccati equation.

From $M = \Psi M$ we get $\|M\|_{C^{0}}\le R$, and the differential equation yields
\[
\Bigl\|\frac{\mathrm{d}M}{\mathrm{d}t}\Bigr\|_{C^{0}} \le \|N\|_{C^{0}} + \frac12 R^{2} = \frac{R}{\ell}
\]
by the choice of $R$.  Thus $\|M\|_{C^{1}} \le R + \frac{R}{\ell} = \bigl(1+\frac1\ell\bigr)R$.
Using $1-\sqrt{1-x}\le x$ for $x\in[0,1)$, we obtain $R\le 2\ell\|N\|_{C^{0}}$, and therefore
\[
\|M\|_{C^{0}} \le 2\ell\|N\|_{C^{0}}, \qquad
\|M\|_{C^{1}} \le 2(\ell+1)\|N\|_{C^{0}} .
\]
This proves (i).

\noindent(ii) From the integral formulation and the bound $\|M_i\|_{C^0}\le 2\ell\delta$,
\[
\begin{aligned}
\|M_1-M_2\|_{C^0}
&\le \ell\|N_1-N_2\|_{C^0} + \frac{\ell}{2}(4\ell\delta)\|M_1-M_2\|_{C^0} \\
&= \ell\|N_1-N_2\|_{C^0} + 2\ell^2\delta\,\|M_1-M_2\|_{C^0}.
\end{aligned}
\]
Hence
\[
\|M_1-M_2\|_{C^0} \le \frac{\ell}{1-2\ell^2\delta}\,\|N_1-N_2\|_{C^0}.
\]
For the derivatives, $\frac{\mathrm{d}M_i}{\mathrm{d}t} = N_i - \frac12 M_i^2$, so using the same bound on $\|M_i\|_{C^0}$,
\[
\begin{aligned}
\Bigl\|\frac{\mathrm{d}M_1}{\mathrm{d}t}-\frac{\mathrm{d}M_2}{\mathrm{d}t}\Bigr\|_{C^0}
&\le \|N_1-N_2\|_{C^0} + \frac12(4\ell\delta)\|M_1-M_2\|_{C^0} \\
&\le \|N_1-N_2\|_{C^0} + 2\ell\delta\cdot\frac{\ell}{1-2\ell^2\delta}\|N_1-N_2\|_{C^0} \\
&= \frac{1}{1-2\ell^2\delta}\,\|N_1-N_2\|_{C^0}.
\end{aligned}
\]
Adding the two estimates gives the required $C^1$ bound.
\end{proof}

\begin{proof}[\textbf{Proof of Theorem~\ref{thm:A} (no-swirl case)}]
The no-swirl case is proved by applying Lemma~\ref{lm:Riccati} with $t_0=0$ and $\ell=1$ to the Riccati equation \eqref{eqn:Riccati}.

\noindent
\underline{Step 1.} \emph{Estimate of $G+2\mcC$.}
For a swirl-free force $\BF = \rho^{-3}(X(\varphi)\Be_\rho + Y(\varphi)\Be_\varphi)$,
\eqref{eqn:sphericalcurl} gives $\nabla\times\BF = -\rho^{-4}(X'+2Y)\Be_\theta$,
and therefore
\[
\|(\curl\BF)^{\tan}\|_{C^0(\mathbb{S}^2\cap\mathbb{R}^3_+)} = \|X'+2Y\|_{C^0[0,\frac{\pi}{2}]}.
\]
Recall that $G(\vp)=X(\vp)+2\int_0^\vp Y(\tau)\, \rmd\tau$ and $\mcC=-\int_0^{\frac{\pi}{2}}G(\tau)\sin\tau \cos\tau \, \rmd\tau$ from \eqref{eqn:constC}.
A direct calculation shows
\begin{align}\label{eqn:Gand2C}
G(\vp)+2\mcC 
= 2\int_0^{\frac{\pi}{2}} \bigl(G(\vp)-G(\tau)\bigr)\sin\tau \cos\tau\, \rmd\tau,
\end{align}
and consequently,
\begin{align*}
\|G+2\mcC\|_{C^0[0,\frac{\pi}{2}]} 
\le \sup_{0 \le \vp,\tau \le \frac{\pi}{2}} |G(\vp)-G(\tau)|
\le \frac{\pi}{2} \|X'+2Y\|_{ C^0[0, \frac{\pi}{2}] }.
\end{align*}
Throughout this proof we keep the coefficients $\pi/2$ explicit because they encode the interval length; this will be convenient in Section~\ref{sec:Discussions} when we replace $[0,\pi/2]$ by $[0,\varphi_0]$ and need to track the angular dependence.

\noindent
\underline{Step 2.} \emph{Estimate of $N$.}
Recall that
\[
N(t)=\frac{-1}{(1-t^2)^2} \int_t^1 J(s) \, \rmd s,
\]
where $J(t)=\int_0^{\vp} (G(\tau) +2\mcC) \sin \tau \, \rmd\tau$ with $t=\cos\vp$.
Using $J(1)=0$, the numerator of $N$ becomes
\begin{equation*}
\begin{aligned}
- \int_t^1 J(s)\, \mathrm{d}s
&= t J(t) + \int_t^1 s \,\frac{\mathrm{d}J}{\mathrm{d}t}(s) \, \mathrm{d}s
= \int_t^1 (s-t)\, \frac{\mathrm{d}J}{\mathrm{d}t}(s)\, \mathrm{d}s \\
&= (1-t)^2 \int_0^1 \frac{\mathrm{d}J}{\mathrm{d}t}\bigl( 1-\theta(1-t) \bigr)(1-\theta) \, \mathrm{d}\theta .
\end{aligned}
\end{equation*}
Therefore
\begin{align}\label{eqn:NandJprime}
N(t) = \frac{1}{(1+t)^2}\int_0^1 \frac{\mathrm{d}J}{\mathrm{d}t}\bigl( 1-\theta(1-t)\bigr)(1-\theta) \, \rmd\theta.
\end{align}
Because $\frac{\mathrm{d}J}{\mathrm{d}t}(t) = -(G(\vp)+2\mcC)$ with $t=\cos\vp$, we have $\bigl\|\frac{\mathrm{d}J}{\mathrm{d}t}\bigr\|_{C^0[0,1]} = \|G+2\mcC\|_{C^0[0,\frac{\pi}2]}$, and
\[
\|N\|_{C^0[0,1]} \le \frac12 \|G+2\mcC\|_{C^0[0,\frac{\pi}{2}]} \le \frac{\pi}{4} \|X'+2Y\|_{ C^0[0, \frac{\pi}{2}] }.
\]

\noindent
\underline{Step 3.} \emph{Existence and uniqueness.}
Set $\varepsilon_1 = \frac{2}{\pi}$. If $\|X'+2Y\|_{C^0[0, \frac{\pi}{2}]} < \varepsilon_1$, then $\|N\|_{C^0[0, 1]} < 1/2$; by Lemma~\ref{lm:Riccati}(i) there exists a unique $M\in C^1[0,1]$. 
Defining $g(\vp)=M(\cos\vp)\sin\vp$, we obtain $g\in C^1[0,\pi/2]$.
To see that $g(\vp)\cot\vp$ is continuous we use $g(0)=0$:
\[
g(\vp)\cot\vp = \cot\vp \int_0^{\vp} g'(\tau)\,\rmd\tau
= \vp\cot\vp \int_0^1 g'(\vp s)\,\rmd s.
\]
Because $|\vp\cot\vp| \le \frac{\pi}{2}$ on $[0,\frac{\pi}{2}]$, the product $g(\varphi)\cot\varphi$ is continuous. Setting $f = -(g' + g\cot\varphi)$ (see \eqref{eqn:halfD2}$_3$) yields $f\in C^0[0,\pi/2]$.  The pressure $P$ is recovered from \eqref{eqn:eqnforP}.  With $h\equiv0$, the quartet $(f,g,h,P)$ satisfies the reduced system \eqref{eqn:halfD1} in the distributional sense. 
Consequently,
\[
\Bu = \frac{f(\varphi)}{\rho}\Be_\rho + \frac{g(\varphi)}{\rho}\Be_\varphi \in C^0(\R^3_+),
\qquad
p = \frac{P(\varphi)}{\rho^2}
\]
fulfil the Navier--Stokes equations \eqref{eqn:NSEhfD} in the sense of distributions.
Uniqueness of $\Bu$ is inherited from the uniqueness of the Riccati equation \eqref{eqn:Riccati}.
Since the angular profiles are $C^1$ and the equations are satisfied away from the origin, standard local Stokes regularity yields \[\n\Bu\in C^{2,\alpha}_{\mathrm{loc}}(\overline{\mathbb R^3_+}\setminus\{0\}),\qquad p\in C^{1,\alpha}_{\mathrm{loc}}(\overline{\mathbb R^3_+}\setminus\{0\})\n\] for every $\alpha\in(0,1)$. Hence all derivatives appearing in the reduction to the ODE system are classical away from the origin.
\end{proof}

On the angular interval $[0,\varphi_0]$ corresponding to a general cone,
the bound $|\varphi\cot\varphi|\le \pi/2$ is replaced by
$|\varphi\cot\varphi|\le M(\varphi_0)$, where $M(\varphi_0)$ is finite for $\varphi_0\in(0,\pi)$,
remains bounded as $\varphi_0\to0$, and blows up only as $\varphi_0\to\pi$, similarly to the constant
$C(\varphi_0)$ in Lemma~\ref{lm:ODEEst}.


\subsection{Blow-up criterion for the swirl-free external force}\label{sec:blowup}

This subsection is devoted to the proof of Theorem~\ref{thm:B}. Recall that in Section~\ref{sec:reductNo-swirl} we derived the Riccati equation
\begin{equation*}
\left\{
\begin{aligned}
\frac{\mathrm{d}M}{\mathrm{d}t} + \frac12 M^{2} &= N, \quad t \in [0,1],\\
M(0)&=0,
\end{aligned}
\right.
\end{equation*}
where
\[
N(t)=-\frac{1}{(1-t^{2})^{2}}\int_{t}^{1} J(s)\,\mathrm{d}s,\qquad
J(t)=\int_{0}^{\varphi} \bigl(G(\tau)+2\mcC\bigr)\sin\tau\,\mathrm{d}\tau,
\]
with $\mcC=-\int_{0}^{\frac{\pi}{2}} G(\tau) \sin\tau \cos\tau \, \mathrm{d}\tau$ and $t=\cos\varphi$.
Replacing $\BF$ by $\lambda\BF$ changes $G$ to $\lambda G$ and $N$ to $\lambda N$, yielding
\begin{equation}\label{eqn:sclRiccati}
\left\{
\begin{aligned}
\frac{\mathrm{d}M}{\mathrm{d}t} + \frac12 M^{2} &= \lambda N, \quad t \in [0,1],\\
M(0)&=0.
\end{aligned}
\right.
\end{equation}


To discuss solutions beyond $t=1$, we extend $N$ constantly to $(-1,2)$ by $N(t)=N(0)$ for $t\in(-1,0]$ and $N(t)=N(1)$ for $t\in[1,2)$.


\begin{theorem}\label{thm:Auxiliary}
Let $N\in C^{0}[0,1]$ and consider the initial value problem \eqref{eqn:sclRiccati} on $(-1,2)$.
For each $\lambda$, let $M_{\lambda}$ be the unique solution of \eqref{eqn:sclRiccati} on $(-1,2)$ and let $[0,\omega_{\lambda})$ be the maximal right interval of existence of the solution. Define
\[
\Lambda:=\{\lambda\in\mathbb{R}\mid\omega_{\lambda}>1\}.
\]
%$\Lambda:=\{\lambda\in\mathbb{R}\mid\omega_{\lambda}>1\}$.
Then $\Lambda$ is an open interval containing $0$; write $\Lambda=(\underline{\lambda},\overline{\lambda})$.
Moreover, the following classification holds.
\begin{enumerate}[(i)]
\item If $N\le0$ on $(0,1)$ and $N\not\equiv0$, then $\underline{\lambda}=-\infty$ and $\overline{\lambda}\in(0,+\infty)$;
\item If $N\ge0$ on $(0,1)$ and $N\not\equiv0$, then $\underline{\lambda}\in(-\infty,0)$ and $\overline{\lambda}=+\infty$;
\item If $N$ changes sign on $(0,1)$, then $\underline{\lambda}\in(-\infty,0)$ and $\overline{\lambda}\in(0,+\infty)$;
\item If $N\equiv0$ on $(0,1)$, then $\underline{\lambda}=-\infty$ and $\overline{\lambda}=+\infty$.
\end{enumerate}
If $\underline{\lambda}$ (resp.\ $\overline{\lambda}$) is finite, then $\omega_{\underline{\lambda}}=1$ (resp.\ $\omega_{\overline{\lambda}}=1$); in other words, the corresponding solution exists on $[0,1)$ and blows up exactly at $t=1$.
\end{theorem}

We prove the theorem with the help of three lemmas.

\begin{lemma}\label{lm:lm1}
Consider the initial value problem \eqref{eqn:sclRiccati}.
If $N$ is negative somewhere in $(0,1)$, then there exists $\lambda^{*}>0$ such that $\omega_{\lambda}\le 1$ for all $\lambda\ge\lambda^{*}$.
Similarly, if $N$ is positive somewhere in $(0,1)$, then there exists $\lambda_{*}<0$ such that $\omega_{\lambda}\le 1$ for all $\lambda\le\lambda_{*}$.
\end{lemma}

\begin{proof}
We prove the first statement; the second follows by symmetry.
Choose $t_{1}\in(0,1)$ and $C>0$ with $N(t_{1})=-2C$. By continuity of $N$, there exists $\delta_{1}>0$ such that $t_{1}-\delta_{1}>0$ and $N(t)<-C$ on $(t_{1}-\delta_{1},t_{1})$.
If the solution does not exist up to $t_{1}-\delta_{1}$, then $\omega_{\lambda}\le t_{1}-\delta_{1}<1$ and the claim holds trivially.
Hence we may assume $\omega_{\lambda}>t_{1}-\delta_{1}$, so that $M(t_{1}-\delta_{1})$ is finite. On the interval $(t_{1}-\delta_{1}, t_{1})$, for any $\lambda>0$ we have
\[
\frac{\mathrm{d}M}{\mathrm{d}t} \le -\Bigl(\frac{1}{2}M^{2}+\lambda C\Bigr),
\]
which gives
\[
\frac{\mathrm{d}}{\mathrm{d}t}\Bigl( \arctan \Bigl(\frac{M(t)}{\sqrt{2\lambda C}}\Bigr) \Bigr)
\le -\sqrt{\frac{\lambda C}{2}}.
\]
Integrating from $t_{1}-\delta_{1}$ to $t$ yields
\[
\arctan \Bigl(\frac{M(t)}{\sqrt{2\lambda C}}\Bigr)
\le -\sqrt{\frac{\lambda C}{2}}\bigl(t-\tilde{t}_{1}\bigr),
\]
where $\tilde{t}_{1} = t_{1}-\delta_{1}+\frac{2}{\sqrt{2\lambda C}}\arctan\bigl(\frac{M(t_{1}-\delta_{1})}{\sqrt{2\lambda C}}\bigr)$.
Hence
\[
M(t) \le -\sqrt{2\lambda C}\, \tan\Bigl(\sqrt{\frac{\lambda C}{2}}(t-\tilde{t}_{1})\Bigr).
\]
Thus $M$ blows down to $-\infty$ no later than $t = \frac{\pi}{\sqrt{2\lambda C}} + \tilde{t}_{1}$.
If $\lambda \ge \lambda^{*} := \frac{2}{C}\bigl(\frac{\pi}{\delta_{1}}\bigr)^{2}$, then
\[
\frac{\pi}{\sqrt{2\lambda C}} + \tilde{t}_{1} 
\le \frac{2\pi}{\sqrt{2\lambda C}} + t_{1} - \delta_{1} \le t_{1} < 1,
\]
so $\omega_{\lambda} \le t_{1} < 1$. This completes the proof.
\end{proof}

\begin{lemma}\label{lm:open}
The set $\Lambda = \{\lambda \in \mathbb{R} \mid \omega_{\lambda} > 1\}$ is an open interval containing $0$.
\end{lemma}

\begin{proof}
\underline{Step 1.} \emph{Openness.}
If $\lambda=0$, equation \eqref{eqn:sclRiccati} becomes $\frac{\mathrm{d}M}{\mathrm{d}t} + \frac12 M^{2} = 0$ with $M(0)=0$, whose unique solution is $M_{0}\equiv0$. By the constant extension of $N$, this solution exists on $(-1,2)$; hence $\omega_{0} = 2 > 1$, so $0\in\Lambda$ and $\Lambda \not =\emptyset$.
Openness follows from the lower semicontinuity of $\omega_{\lambda}$ with respect to $\lambda$ (see \cite[Chapter~V, Theorem~2.1]{Hartman02}). Indeed, for $\lambda_{1}\in\Lambda$ we have $\omega_{\lambda_{1}}>1$; choose $\delta_2>0$ such that $\omega_{\lambda_{1}}>1+\delta_2$. By lower semicontinuity there exists $\varepsilon>0$ such that
$\omega_{\lambda} \ge \omega_{\lambda_{1}} - \frac{\delta_2}{2} > 1 + \frac{\delta_2}{2} > 1$ whenever $|\lambda-\lambda_{1}|<\varepsilon$. Thus $(\lambda_{1}-\varepsilon,\lambda_{1}+\varepsilon)\subset\Lambda$, proving openness.

\underline{Step 2.} \emph{Interval property.}
Since $\Lambda$ is open and contains $0$, it is a union of disjoint open intervals; let $(\underline{\lambda},\overline{\lambda})$ be the connected component containing $0$. We show that actually $\Lambda = (\underline{\lambda},\overline{\lambda})$. Denote $C_{1} := \max\limits_{t\in[0,1]}|N(t)|$.

\underline{Case 1.} \emph{$N \le 0$ on $(0,1)$ and $N \not\equiv 0$.}
Lemma~\ref{lm:lm1} forces $0<\overline{\lambda}<+\infty$ and the openness of $\Lambda$ gives $\omega_{\overline{\lambda}}\le 1$.
We prove that $\underline{\lambda}=-\infty$ and that $\omega_{\lambda}\le 1$ for all $\lambda>\overline{\lambda}$.
Since $\omega_{\lambda}>1$ for $\lambda\in[0,\overline{\lambda})$, it suffices to prove the same inequality for all $\lambda<0$; this will force $\underline{\lambda}=-\infty$.
To this end, fix any $\lambda<0$ and choose a $\lambda_2 \in (0,\overline{\lambda})$; then $\omega_{\lambda_2}>1$.
Since $N\le0$ and $\lambda<\lambda_{2}$, we have
\[
\lambda_{2} N - \frac12 M^{2} \le \lambda N - \frac12 M^{2} \le |\lambda| C_{1} - \frac12 M^{2}.
\]
By standard comparison theorems (see \cite[Chapter~III, Corollary~4.2]{Hartman02}),
\begin{align}\label{eqn:lm2Comp}
M_{\lambda_{2}}(t) \le M_{\lambda}(t) \le \sqrt{2|\lambda| C_{1}}.
\end{align}
Thus $\omega_{\lambda} \ge \omega_{\lambda_{2}} > 1$ for every $\lambda<0$, forcing $\underline{\lambda}=-\infty$.
Moreover, the upper bound in \eqref{eqn:lm2Comp} shows that any finite-time singularity is a blow-down to $-\infty$.
Since $\omega_{\overline{\lambda}} \le 1$, $M_{\overline{\lambda}}$ blows down no later than $t=1$.

Now take any $\lambda > \overline{\lambda}>0$ and define $v(t) = \lambda\frac{M_{\overline{\lambda}}(t)}{\overline{\lambda}}$.
A direct computation gives
\[
\frac{\mathrm{d}v}{\mathrm{d}t} = \lambda N - \frac12 \frac{\overline{\lambda}}{\lambda} v^{2}
\ge \lambda N - \frac12 v^{2},
\]
where we used $0<\overline{\lambda}/\lambda<1$. Since $M_{\lambda}(0)=v(0)=0$, comparison gives $M_{\lambda}\le v$. Because $v$ blows down no later than $t=1$, we get $\omega_{\lambda}\le1$ for all $\lambda>\overline{\lambda}$.

\underline{Case 2.} \emph{$N \ge 0$ on $(0,1)$ and $N \not\equiv 0$.}
By symmetry, Lemma~\ref{lm:lm1} gives $-\infty<\underline{\lambda}<0$, openness yields $\omega_{\underline{\lambda}}\le1$, and the same comparison argument shows $\overline{\lambda}=+\infty$ and $\omega_{\lambda}\le1$ for all $\lambda<\underline{\lambda}$. We omit the details.

\underline{Case 3.} \emph{$N$ changes sign on $(0,1)$.}
Lemma~\ref{lm:lm1} implies $-\infty<\underline{\lambda}<0<\overline{\lambda}<+\infty$, together with $\omega_{\underline{\lambda}}\le1$ and $\omega_{\overline{\lambda}}\le1$.
We must show $\omega_{\lambda}\le1$ for $\lambda\notin[\underline{\lambda},\overline{\lambda}]$.
From \eqref{eqn:lm2Comp} we know that any blow-up is a blow-down to $-\infty$, and $M_{\underline{\lambda}},M_{\overline{\lambda}}$ blow down no later than $t=1$.
For $\lambda > \overline{\lambda}$, the scaling $v(t)=\lambda\frac{M_{\overline{\lambda}}(t)}{\overline{\lambda}}$ again satisfies $\frac{\mathrm{d}v}{\mathrm{d}t}\ge \lambda N - \frac12 v^{2}$, hence $M_{\lambda}\le v$ and $\omega_{\lambda}\le1$.
For $\lambda < \underline{\lambda}$, setting $v(t)=\lambda\frac{M_{\underline{\lambda}}(t)}{\underline{\lambda}}$ yields $\frac{\mathrm{d}v}{\mathrm{d}t}\ge \lambda N - \frac12 v^{2}$, since $0<\underline{\lambda}/\lambda<1$; therefore $M_{\lambda}\le v$ and $\omega_{\lambda}\le1$.
\end{proof}

\begin{lemma}\label{lm:exactblow}
If $\underline{\lambda}$ (resp.\ $\overline{\lambda}$) is finite, then $\omega_{\underline{\lambda}}=1$ (resp.\ $\omega_{\overline{\lambda}}=1$).
\end{lemma}

\begin{proof}
We prove $\omega_{\overline{\lambda}}=1$ by contradiction; the case for $\underline{\lambda}$ is analogous.
Suppose $\omega_{\overline{\lambda}}<1$. Choose $t_{R}$ with $t_{R}<\omega_{\overline{\lambda}}<1$ and a number $R>0$ to be specified later such that $M_{\overline{\lambda}}(t_{R})<-2R$.
By continuous dependence on the parameter (see \cite[Chapter~V, Theorem~2.1]{Hartman02}), there exists $\lambda_{3}\in(\underline{\lambda},\overline{\lambda})$ close to $\overline{\lambda}$ with $M_{\lambda_{3}}(t_{R}) < -R$.
Let $C_{1} := \max\limits_{t\in[0,1]}|N(t)|$. Then
\[
\frac{\mathrm{d}M_{\lambda_{3}}}{\mathrm{d}t} = \lambda_{3} N - \frac12 M_{\lambda_{3}}^{2} \le \overline{\lambda} C_{1} - \frac12 M_{\lambda_{3}}^{2}.
\]
Consider the auxiliary problem
\[
\frac{\mathrm{d}v}{\mathrm{d}t} = \overline{\lambda} C_{1} - \frac12 v^{2},\qquad v(t_{R}) = -R.
\]
Choose $R > \sqrt{4\overline{\lambda} C_{1}}$. Then $\frac{\mathrm{d}v}{\mathrm{d}t}(t_{R})<0$, and one checks that $v(t)\le -R$, $\frac{\mathrm{d}v}{\mathrm{d}t}(t)<0$ for all $t\ge t_{R}$ as long as the solution exists.
Moreover, for $v\le -R$ we have $\frac14 v^{2} \ge \overline{\lambda} C_{1}$, hence
\[
\frac{\mathrm{d}v}{\mathrm{d}t} = \overline{\lambda} C_{1} - \frac12 v^{2} \le \frac14 v^{2} - \frac12 v^{2} = -\frac14 v^{2}.
\]
Integrating from $t_{R}$ to $t$ gives
\[
-\frac{1}{v(t)} \le \frac{1}{R} - \frac14 (t - t_{R}),
\]
%$-\frac{1}{v(t)} \le \frac{1}{R} - \frac14 (t - t_{R})$, 
and because $v(t)<0$, the right-hand side must stay positive, which forces
$t < t_{R} + \frac{4}{R}$.
Thus $v$ blows down to $-\infty$ no later than $t_{R} + 4/R$.
Choosing $R > \max\bigl\{ \sqrt{4\overline{\lambda} C_{1}},\; \frac{4}{1-\omega_{\overline{\lambda}}} \bigr\}$ guarantees $t_{R} + 4/R < 1$.
By comparison, $M_{\lambda_{3}}(t)\le v(t)$ on their common interval, so $M_{\lambda_{3}}$ also blows down before $t=1$, contradicting $\omega_{\lambda_{3}}>1$. Hence $\omega_{\overline{\lambda}}=1$.
\end{proof}

\begin{proof}[\textbf{Proof of Theorem~\ref{thm:B}}]
Theorem~\ref{thm:B} follows directly from Theorem~\ref{thm:Auxiliary}. Indeed, the Navier--Stokes system with force $\lambda\BF$ reduces to the Riccati equation \eqref{eqn:sclRiccati}. By Theorem~\ref{thm:Auxiliary}, the set of admissible $\lambda$ is an open interval $(\underline{\lambda},\overline{\lambda})$ characterized by the sign of $N$ on $(0,1)$.
It remains to verify that the discriminant $\mathcal{N}$ and the function $N$ have the same sign.
Since $N(t)=-\frac{1}{(1-t^{2})^{2}}\int_{t}^{1} J(s)\,\mathrm{d}s$, the sign of $N$ is determined by $-\int_{t}^{1} J(s)\,\mathrm{d}s$.
Using the change of variable $s=\cos\tau$, we compute
\begin{equation}\label{eqn:PfThmB}
\begin{aligned}
-\int_{t}^{1} J(s)\, \mathrm{d}s
&= -\int_{\varphi}^{0} \big( \int_0^{\tau} (G(\sigma) + 2\mcC) \sin \sigma \, \mathrm{d}\sigma \big) \, \mathrm{d}(\cos\tau)\\
&= \cos \varphi \int_0^{\vp} (G(\tau) + 2\mcC) \sin \tau \, \mathrm{d}\tau  - \int_{0}^{\varphi} (G(\tau)+2\mcC)\sin \tau \cos \tau \, \mathrm{d}\tau\\
&= \int_{0}^{\varphi} (G(\tau)+2\mcC) \sin \tau (\cos \varphi - \cos \tau) \, \mathrm{d}\tau = \mathcal{N}(\varphi).
\end{aligned}
\end{equation}
Thus $\operatorname{sgn} N(t) = \operatorname{sgn} \mathcal{N}(\varphi)$, and the classification in Theorem~\ref{thm:Auxiliary} translates exactly into the four cases listed in Theorem~\ref{thm:B}.
\end{proof}

\begin{remark}[Blow-up location]
If $\underline{\lambda}$ (resp.\ $\overline{\lambda}$) is finite, the equality $\omega_{\underline{\lambda}}=1$ (resp.\ $\omega_{\overline{\lambda}}=1$) means that the solution $M(t)$ of the Riccati equation exists on $[0,1)$ and blows down as $t\to 1^{-}$. In terms of $\varphi$, this corresponds to a solution that is well defined on $(0,\pi/2]$ but blows down exactly at the north pole $\varphi=0$.
\end{remark}


%% FIXED: reorganized floating text into a single coherent remark
\begin{remark}[Structure of the discriminant]
Introduce the weight functions
\[
\mathcal{W}_1(s)=2\sin s\cos s,\qquad
\mathcal{W}_{2,\varphi}(\tau)=\sin\tau(\cos\varphi-\cos\tau).
\]
For a swirl-free force $\BF=\rho^{-3}(X\Be_\rho+Y\Be_\varphi)$, the curl reduces to
$\nabla\times\BF = -\rho^{-4}(X'+2Y)\Be_\theta$.
From \eqref{eqn:disfcn} and the definition of $\mcC$ one obtains
\[
\mathcal{N}(\varphi)
= \int_{0}^{\varphi} \mathcal{W}_{2,\varphi}(\tau)\bigl(G(\tau)+2\mcC\bigr)\,\mathrm{d}\tau
= \int_{0}^{\varphi} \mathcal{W}_{2,\varphi}(\tau)
   \Biggl[ \int_{0}^{\pi/2} \mathcal{W}_1(s)
      \biggl( \int_{s}^{\tau} (X'+2Y)(t)\,\mathrm{d}t \biggr)
      \,\mathrm{d}s \Biggr] \mathrm{d}\tau .
\]
Thus $\mathcal{N}$ is precisely the tangential curl $X'+2Y$ integrated first along the meridian
and then against the two spherical-cap weight functions $\mathcal{W}_1$ and $\mathcal{W}_{2,\varphi}$.
Consequently, the sign of $\mathcal{N}(\varphi)$ is completely determined by the distribution of the
tangential curl.

The discriminant $\mathcal{N}$ also satisfies
\begin{align}\label{eqn:discriminantN}
\mathcal{L}[\mathcal{N}] = -(G+2\mcC)\sin^{2}\varphi,\qquad
\mathcal{N}(0)=\mathcal{N}(\tfrac{\pi}{2})=0,
\end{align}
where $\mathcal{L}= \frac{\mathrm{d}^{2}}{\mathrm{d}\varphi^{2}}-\cot\varphi\frac{\mathrm{d}}{\mathrm{d}\varphi}
= \sin\varphi\frac{\mathrm{d}}{\mathrm{d}\varphi}\Bigl(\frac{1}{\sin\varphi}\frac{\mathrm{d}}{\mathrm{d}\varphi}\Bigr)$.
This is precisely the angular part of the Stokes stream-function operator
$E^{2}= \frac{\partial^{2}}{\partial\rho^{2}} + \frac{\sin\varphi}{\rho^{2}}\frac{\partial}{\partial\varphi}\!\bigl(\frac{1}{\sin\varphi}\frac{\partial}{\partial\varphi}\bigr)$
(see, e.g., \cite[eq.~(4-7.18)]{HB65}) under the spherical coordinates.
Solving the Dirichlet problem for $\mathcal{L}$ gives the integral representation of $\mathcal{N}$, thus
$\mathcal{W}_{2,\varphi}(\tau)=\sin\tau(\cos\varphi-\cos\tau)$ is exactly the Green's function of $\mathcal{L}$
multiplied by the source weight $-\sin^{2}\tau$.

%Notice that evaluating \eqref{eqn:PfThmB} at $t=0$ (i.\,e.\ $\varphi=\pi/2$) gives
%\[
%\int_{0}^{1} J(s)\,\mathrm{d}s = \int_{0}^{\pi/2} \bigl(G(\varphi)+2\mcC\bigr)\sin\varphi\cos\varphi\,\mathrm{d}\varphi .
%\]
%Combined with the compatibility condition $\int_{0}^{1} J(s)\,\mathrm{d}s = 0$ and the positivity of
%$\sin\varphi\cos\varphi$ on $(0,\pi/2)$, this forces $G+2\mcC$ to change sign in $(0,\pi/2)$.  Since $G+2\mcC$ must change sign, the source term in \eqref{eqn:discriminantN} is sign-changing; the maximum principle for $\mathcal{L}$ does not determine the sign of $\mathcal{N}$, and three distinct cases arise, as illustrated in Example \ref{example:N}.
\end{remark}








%Introduce the weight functions
%\[
%\omega_1(s)=2\sin s\cos s,\qquad
%\omega_{2,\varphi}(\tau)=\sin\tau(\cos\varphi-\cos\tau).
%\]
%For a swirl-free force $\BF=\rho^{-3}(X\Be_\rho+Y\Be_\varphi)$, the curl reduces to
%$\nabla\times\BF = -\rho^{-4}(X'+2Y)\Be_\theta$.
%From \eqref{eqn:disfcn} and the definition of $\mathcal{C}$ one obtains
%\[
%\mathcal{N}(\varphi)
%= \int_{0}^{\varphi} \omega_{2,\varphi}(\tau)\bigl(G(\tau)+2\mathcal{C}\bigr)\,\mathrm{d}\tau
%= \int_{0}^{\varphi} \omega_{2,\varphi}(\tau)
 %  \Biggl[ \int_{0}^{\pi/2} \omega_1(s)
  %    \biggl( \int_{s}^{\tau} (X'+2Y)(t)\,\mathrm{d}t \biggr)
   %   \,\mathrm{d}s \Biggr] \mathrm{d}\tau .
%\]
%Thus $\mathcal{N}$ is precisely the tangential curl %$X'+2Y$ integrated first along the meridian
%and then against the two spherical-cap weight functions $\omega_1$ and $\omega_{2,\varphi}$.
%Consequently, the sign of $\mathcal{N}(\varphi)$ is completely determined by the distribution of the
%tangential curl.


%\begin{remark}[Origin of the weight $\omega_{2,\varphi}$]
%The discriminant $\mathcal{N}$ satisfies
%\begin{align}\label{eqn:discriminantN}
%\mathcal{L}[\mathcal{N}] = -(G+2\mathcal{C})\sin^{2}\varphi,\qquad
%\mathcal{N}(0)=\mathcal{N}(\tfrac{\pi}{2})=0,
%\end{align}
%where $\mathcal{L}= \frac{\mathrm{d}^{2}}{\mathrm{d}\varphi^{2}}-\cot\varphi\frac{\mathrm{d}}{\mathrm{d}\varphi}
%= \sin\varphi\frac{\mathrm{d}}{\mathrm{d}\varphi}\Bigl(\frac{1}{\sin\varphi}\frac{\mathrm{d}}{\mathrm{d}\varphi}\Bigr)$.
%This is precisely the angular part of the Stokes stream-function operator
%$E^{2}= \frac{\partial^{2}}{\partial\rho^{2}} + \frac{\sin\varphi}{\rho^{2}}\frac{\partial}{\partial\varphi}\!\bigl(\frac{1}{\sin\varphi}\frac{\partial}{\partial\varphi}\bigr)$
%(see, e.g., \cite[Eq.~(4-7.18)]{HB65}) under the self-similar ansatz.
%Solving the Dirichlet problem for $\mathcal{L}$ gives the integral representation of $\mathcal{N}$, thus
%$\omega_{2,\varphi}(\tau)=\sin\tau(\cos\varphi-\cos\tau)$ is exactly the Green's function of $\mathcal{L}$
%multiplied by the source weight $-\sin^{2}\tau$.
%\end{remark}


Notice that evaluating \eqref{eqn:PfThmB} at $t=0$ (i.\,e.\ $\varphi=\pi/2$) gives
\[
\int_{0}^{1} J(s)\,\mathrm{d}s = \int_{0}^{\pi/2} \bigl(G(\varphi)+2\mathcal{C}\bigr)\sin\varphi\cos\varphi\,\mathrm{d}\varphi .
\]
Combined with the compatibility condition $\int_{0}^{1} J(s)\,\mathrm{d}s = 0$ and the positivity of
$\sin\varphi\cos\varphi$ on $(0,\pi/2)$, this forces $G+2\mcC$ to change sign in $(0,\pi/2)$.  Since $G+2\mathcal{C}$ must change sign, the source term in \eqref{eqn:discriminantN} is sign-changing; the maximum principle for $\mathcal{L}$ does not determine the sign of $\mathcal{N}$, and three distinct cases arise, as illustrated in Example \ref{example:N}.






\section{Analysis of solutions for a swirling external force}\label{sec:withswirl}


\subsection{Reduction for external forces with swirl}\label{sec:reduct-swirl}

We now consider the general case where the external force $\BF$ possesses swirl. The system reads
\begin{equation}\label{eqn:hDgeneral1}
\left\{
\begin{aligned}
   f^{\prime\prime} + f^{\prime} \cot \varphi &= gf^{\prime} -(f^2 + g^2 + h^2) - 2P - X, \\
   f^{\prime} &= gg^{\prime} - h^2 \cot \varphi + P^{\prime} - Y, \\ 
   (h^\prime + h\cot \varphi)^{\prime} &= g(h^\prime + h\cot \varphi) - Z,\\
   f+ g^\prime + g\cot \varphi &=0,
\end{aligned}
\right.
\end{equation}
together with the boundary conditions
\begin{align*}%\label{eqn:hDgeneralBC}
f^\prime(0) = g(0) = h(0) = 0, \quad
  f\Bigl(\frac{\pi}{2}\Bigr) = g\Bigl(\frac{\pi}{2}\Bigr) = h\Bigl(\frac{\pi}{2}\Bigr) = 0.
\end{align*}

From \eqref{eqn:hDgeneral1}$_2$ we obtain 
\[
\frac12 g^{2} + P = f + \int_{0}^{\varphi}\!\bigl( Y(\tau) + h^{2}\cot\tau \bigr)\,\mathrm{d}\tau + \mathcal{C}.
\]
Substituting this into \eqref{eqn:hDgeneral1}$_1$ and using the continuity equation \eqref{eqn:hDgeneral1}$_4$ yields an equation identical to \eqref{eqn:reduct1},
but now with $G$ given by
\begin{equation}\label{eqn:genNotation}
G(\varphi) = X(\varphi) + h^{2}(\varphi) + 2\int_{0}^{\varphi} \bigl( Y(\tau) + h^{2}(\tau)\cot\tau \bigr)\,\mathrm{d}\tau .
\end{equation}

Proceeding now exactly as in Section~\ref{sec:reductNo-swirl} leads to the Riccati equation
\[
\frac{\mathrm{d}M}{\mathrm{d}t} + \frac12 M^{2} = N,\quad t\in[0,1],\qquad M(0)=0,
\]
where $t=\cos\varphi$, $M(t)=\frac{g(\varphi)}{\sin\varphi}$, and
\[
N(t) = -\frac{1}{(1-t^{2})^{2}}\int_{t}^{1} J(s)\,\mathrm{d}s,\quad
J(t) = \int_{0}^{\varphi} \bigl( G(\tau) + 2\mathcal{C} \bigr)\sin\tau\,\mathrm{d}\tau .
\]

The constant $\mathcal{C}$ is fixed by the compatibility condition $\int_{0}^{1} J(s)\,\mathrm{d}s=0$.
Writing this condition explicitly and inserting \eqref{eqn:genNotation} yields
\begin{align}\label{eqn:swirlconstC}
\mathcal{C}
= -\int_{0}^{\pi/2} G(\tau)\sin\tau\cos\tau\,\mathrm{d}\tau 
= -\int_{0}^{\pi/2}\big( X\sin\varphi\cos\varphi + Y\cos^{2}\varphi
+ h^{2}\cot\vp\big)\,\mathrm{d}\varphi ,
\end{align}
where we used the identity
\[
\int_{0}^{\pi/2} h^{2}\sin\varphi\cos\varphi\,\mathrm{d}\varphi
+ 2\int_{0}^{\pi/2}\!\biggl( \int_{0}^{\varphi} h^{2}(\tau)\cot\tau\,\mathrm{d}\tau \biggr)\sin\varphi\cos\varphi\,\mathrm{d}\varphi
= \int_{0}^{\pi/2} h^{2}\cot\varphi\,\mathrm{d}\varphi.
\]

Thus the full system \eqref{eqn:hDgeneral1} reduces to the above Riccati equation coupled with the linear boundary value problem for $h$,
\begin{equation}\label{eqn:swirlh}
(h' + h\cot\varphi)' = g\,(h' + h\cot\varphi) - Z, \quad \vp \in \Bigl[0,\frac{\pi}{2}\Bigr],
\qquad h(0)=h\Bigl(\frac{\pi}{2}\Bigr)=0,
\end{equation}
where $g(\varphi)=M(\cos\varphi)\sin\varphi$.  If $Z\equiv0$, then $h\equiv0$ and the system reduces to the swirl-free case treated in Section~\ref{sec:reductNo-swirl}.


\subsection{Existence and uniqueness for small forces with swirl}\label{sec:ExistSwirl}


We now establish the estimates for the boundary value problem \eqref{eqn:swirlh}.
For later use on solid cones, the lemma is stated and proved on an arbitrary angular interval $[0,\varphi_0]$ with $\varphi_0\in(0,\pi)$. The half-space case is recovered by taking $\varphi_0=\frac{\pi}{2}$.

\begin{lemma}\label{lm:ODEEst}
Consider the boundary value problem
\[
(h' + h\cot \varphi)' = g\,(h' + h\cot \varphi) - Z, \quad \vp \in [0,\varphi_0],\qquad h(0)=h(\varphi_0)=0,
\]
with $\varphi_0\in(0,\pi)$ and $Z\in C^0[0,\varphi_0]$.
\begin{enumerate}[(i)]
\item If $g\in C^0[0,\varphi_0]$, there exists a unique solution $h\in C^1[0,\varphi_0]$, and
\[
\|h\|_{C^1[0,\varphi_0]} \le C(\varphi_0)\Bigl(2\varphi_0 e^{2\varphi_0\|g\|_{C^0[0,\varphi_0]}}(1+\|g\|_{C^0[0,\varphi_0]}) + 1\Bigr)\|Z\|_{C^0[0,\varphi_0]}.
\]
%If $g\in C^1[0,\varphi_0]$, the same estimate holds with $\|g\|_{C^0[0,\varphi_0]}$ replaced by $\|g\|_{C^1[0,\varphi_0]}$.
\item For $g_1,g_2\in C^0[0,\varphi_0]$ and the same $Z$, let $h_1,h_2$ be the corresponding solutions.
Then
\begin{align*}
&\|h_1-h_2\|_{C^1[0,\varphi_0]} \\
&\le C(\varphi_0)\prod_{i=1}^2
\Bigl(2\varphi_0 e^{2\varphi_0\|g_i\|_{C^0[0,\varphi_0]}}(1+\|g_i\|_{C^0[0,\varphi_0]}) + 1\Bigr)
\|Z\|_{C^0[0,\varphi_0]}\|g_1-g_2\|_{C^0[0,\varphi_0]}.    
\end{align*}
%If $g_1,g_2\in C^1$, the same estimate holds with $\|g_i\|_{C^0[0,\varphi_0]}$ replaced by $\|g_i\|_{C^1[0,\varphi_0]}$.
\end{enumerate}
The constant $C(\varphi_0)$ depends only on $\varphi_0$.
\end{lemma}

\begin{proof}
(i) For simplicity we omit the interval $[0,\varphi_0]$ in the norm notation.
Set $H = h' + h\cot\varphi$.  From $H\sin\varphi = (h\sin\varphi)'$ and the boundary conditions,
\[
H' - gH = -Z,\qquad \int_0^{\varphi_0} H(\varphi)\sin\varphi\,\mathrm{d}\varphi = 0.
\]
The linear ODE gives $H(\varphi) = \tilde g(\varphi)(H(0) - \tilde Z(\varphi))$ where
$\tilde g = \exp\int_0^\varphi g$, $\tilde Z = \int_0^\varphi Z\exp(-\int_0^t g)\,\rmd t$.
The integral condition yields
\[
H(0) = \frac{\int_0^{\varphi_0}\tilde g \tilde Z \sin\varphi\,\mathrm{d}\varphi}
            {\int_0^{\varphi_0}\tilde g \sin\varphi\,\mathrm{d}\varphi},
\]
hence $|H(0)| \le \|\tilde Z\|_{C^0} \le \varphi_0 e^{\varphi_0\|g\|_{C^0}}\|Z\|_{C^0}$.
Then
\begin{align}\label{eqn:lm2EstH}
\|H\|_{C^0} \le e^{\varphi_0\|g\|_{C^0}}\bigl(|H(0)|+\|\tilde Z\|_{C^0}\bigr)
          \le 2\varphi_0 e^{2\varphi_0\|g\|_{C^0}}\|Z\|_{C^0}.
\end{align}
From $H' = gH - Z$ we get
\[
\|H'\|_{C^0} \le \|g\|_{C^0}\|H\|_{C^0} + \|Z\|_{C^0}
\le \bigl(2\varphi_0 e^{2\varphi_0\|g\|_{C^0}}\|g\|_{C^0} + 1\bigr)\|Z\|_{C^0},
\]
so that
\[
\|H\|_{C^1} \le \bigl(2\varphi_0 e^{2\varphi_0\|g\|_{C^0}}(1+\|g\|_{C^0}) + 1\bigr)\|Z\|_{C^0}.
\]
The function $h$ is recovered by
\begin{align}\label{eqn:BVPKernel}
h(\varphi) = \frac{1}{\sin\varphi}\int_0^\varphi H(\tau)\sin\tau\,\rmd\tau
= \int_0^1 H(\varphi s)\frac{\varphi\sin(\varphi s)}{\sin\varphi}\,\rmd s.
\end{align}
L'H\^opital's rule and the integral condition give $h(0)=h(\varphi_0)=0$.
The kernel $K(\varphi,s)=\frac{\varphi\sin(\varphi s)}{\sin\varphi}$ is real analytic on $[0,\varphi_0]\times[0,1]$ because $\varphi_0<\pi$. Hence $\|h\|_{C^{1}}\le C(\varphi_{0})\|H\|_{C^{1}}$ for a constant $C(\varphi_{0})$ depending only on $\varphi_{0}$. This proves (i).

\noindent(ii) Let $h_1,h_2$ correspond to $g_1,g_2$.  Set $\tilde h = h_1-h_2$, $\tilde H = \tilde h'+\tilde h\cot\varphi$, and $H_2 = h_2'+h_2\cot\varphi$.  Subtracting the equations,
\[
\tilde H' - g_1\tilde H = (g_1-g_2)H_2,\qquad \tilde h(0)=\tilde h(\varphi_0)=0.
\]
Since $\tilde H$ fulfills the hypotheses of part (i) with $g=g_1$ and $Z$ replaced by $-(g_1-g_2)H_2$, its estimate gives
\[
\|\tilde h\|_{C^1} \le C(\varphi_0)\bigl(2\varphi_0 e^{2\varphi_0\|g_1\|_{C^0}}(1+\|g_1\|_{C^0}) + 1\bigr)
   \|H_2\|_{C^0}\,\|g_1-g_2\|_{C^0}.
\]
Using \eqref{eqn:lm2EstH} for $h_{2}$ we have
\[
\|H_2\|_{C^0} \le \bigl(2\varphi_0 e^{2\varphi_0\|g_2\|_{C^0}}(1+\|g_2\|_{C^0}) + 1\bigr)\|Z\|_{C^0}.
\]
Multiplying the two inequalities gives the desired estimate.
\end{proof}


%\begin{remark}\label{rmk:C_vp_0}
From \eqref{eqn:BVPKernel}, the constant $C(\varphi_{0})$ in Lemma~\ref{lm:ODEEst} is controlled by the $C^{1}$-norm of the kernel
$K(\varphi,s)=\frac{\varphi\sin(\varphi s)}{\sin\varphi}$ on $[0,\varphi_{0}]\times[0,1]$.
It is finite for every $\varphi_{0}\in(0,\pi)$, remains bounded as $\varphi_{0}\to0$,
and blows up as $\varphi_{0}\to\pi$.
%\end{remark}


\begin{proof}[\textbf{Proof of Theorem~\ref{thm:A} (swirl case)}]
\underline{Step 1.} \emph{Estimate of $h$.}
For $\BF = \frac{1}{\rho^3}(X(\varphi)\Be_\rho + Y(\varphi)\Be_\varphi + Z(\varphi)\Be_\theta)$, \eqref{eqn:sphericalcurl} gives
\[
\|(\curl\, \BF)^{\tan}\|_{C^0(\mathbb{S}^2 \cap \R^3_+)} = \bigl\|2Z\Be_{\varphi} - (X'+2Y)\Be_{\theta}\bigr\|_{C^0(\mathbb{S}^2 \cap \R^3_+)},\quad \text{comparable to }2\|Z\|_{C^0[0,\frac{\pi}{2}]} + \|X'+2Y\|_{C^0[0,\frac{\pi}{2}]}.
\]
Fix $g\in C^1[0,\pi/2]$ and $Z\in C^0[0,\pi/2]$. Lemma~\ref{lm:ODEEst}(i) applied to
\begin{equation}\label{eqn:swirlh1}
(h' + h\cot\varphi)' = g(h' + h\cot\varphi) - Z, \quad \vp \in [0,\tfrac{\pi}{2}],\qquad h(0)=h(\tfrac{\pi}{2})=0,
\end{equation}
yields a unique solution $h\in C^1[0,\pi/2]$ with
\begin{equation}\label{eqn:swirlEst1}
\|h\|_{C^1[0,\frac{\pi}{2}]} \le C(\|g\|_{C^1[0,\frac{\pi}{2}]})\,\|Z\|_{C^0[0,\frac{\pi}{2}]},
\end{equation}
where the constant $C(\cdot)$ depends only on $\|g\|_{C^1[0,\frac{\pi}{2}]}$.

\noindent\underline{Step 2.} \emph{Estimate of $N$.}
We now estimate $N$ by merging Steps~1 and~2 of the swirl-free proof with the additional terms coming from the swirl component~$h$.
Define
\[
N(t) = -\frac{1}{(1-t^2)^2}\int_t^1 J(s)\,\mathrm{d}s,
\quad 
J(t) = \int_0^{\varphi} \bigl(G(\tau)+2\mathcal{C}\bigr)\sin\tau\,\mathrm{d}\tau,
\]
with $t=\cos\varphi$ and $\mcC$ from \eqref{eqn:swirlconstC}, and 
\begin{align}\label{eqn:coneG}
G(\varphi) = X(\varphi) + h^2(\varphi) + 2\int_0^{\varphi} \bigl( Y(\tau) + h^2(\tau)\cot\tau \bigr)\,\mathrm{d}\tau.
\end{align}
Exactly as in the swirl-free case,  \eqref{eqn:NandJprime} and \eqref{eqn:Gand2C} give
\[
\|N\|_{C^0[0,1]} \le \frac12 \|G+2\mathcal{C}\|_{C^0[0,\frac{\pi}{2}]}
\le \frac12 \sup_{0\le\tau,\varphi\le\pi/2} |G(\varphi)-G(\tau)|
\le \frac{\pi}{4} \|G'\|_{C^0[0,\frac{\pi}{2}]},
\]
and we continue to keep the factor $\pi/2$ explicit for the extension to a cone in Section~\ref{sec:Discussions}.
Using $h(0)=0$, $\varphi\cot\varphi\le1$ on $[0,\pi/2]$ and the estimate~\eqref{eqn:swirlEst1} for $h$, we obtain
\begin{align*}
|2h(\varphi)h'(\varphi)|
&\le 2\varphi\|h'\|_{C^{0}[0,\frac{\pi}{2}]}^{2}
\le \pi C(\|g\|_{C^{1}[0,\frac{\pi}{2}]})\|Z\|_{C^{0}[0,\frac{\pi}{2}]}^{2},
\end{align*}
and 
\begin{align*}
|2h^{2}(\varphi)\cot\varphi|
&\le 2\varphi^{2}\|h'\|_{C^{0}[0,\frac{\pi}{2}]}^{2}\cot\varphi
\le \pi C(\|g\|_{C^{1}[0,\frac{\pi}{2}]})\|Z\|_{C^{0}[0,\frac{\pi}{2}]}^{2}.
\end{align*}
Hence
\[
\|G'\|_{C^{0}[0,\frac{\pi}{2}]}
\le \|X'+2Y\|_{C^{0}[0,\frac{\pi}{2}]} + \pi C(\|g\|_{C^{1}[0,\frac{\pi}{2}]})\|Z\|_{C^{0}[0,\frac{\pi}{2}]}^{2},
\]
and therefore
\begin{equation}\label{eqn:swirlEstN}
\|N\|_{C^{0}[0,1]}
\le \frac{\pi}{4}\,\|X'+2Y\|_{C^{0}[0,\frac{\pi}{2}]}
+ \frac{\pi^{2}}{4}\,C(\|g\|_{C^{1}[0,\frac{\pi}{2}]})\,\|Z\|_{C^{0}[0,\frac{\pi}{2}]}^{2}.
\end{equation}


\noindent\underline{Step 3.} \emph{Contraction and existence.}
Consider the Riccati problem
\[
\frac{\rmd M}{\rmd t} + \frac12 M^2 = N, \quad t \in [0,1],\qquad M(0)=0.
\]
Set $\mathfrak{Y}_3 := \{ g\in C^1[0,\frac{\pi}{2}] : g(0)=g(\pi/2)=0,\ \|g\|_{C^1[0,\frac{\pi}{2}]}\le 3\}$.
For any $g\in\mathfrak{Y}_{3}$, estimate~\eqref{eqn:swirlEstN} shows that by taking $\|X'+2Y\|_{C^{0}[0,\frac{\pi}{2}]}$ and $\|Z\|_{C^{0}[0,\frac{\pi}{2}]}$ sufficiently small, we can ensure $\|N\|_{C^{0}[0,1]}<\frac14$.
By Lemma~\ref{lm:Riccati}(i) there exists a unique $M\in C^{1}[0,1]$ with $\|M\|_{C^{1}[0,1]}<1$.
Define $Tg(\varphi) := M(\cos\varphi)\sin\varphi$.
The mapping $T$ is obtained through the following steps:
\begin{figure}[htbp]
\centering
\begin{tikzpicture}[>=Stealth, node distance=2cm, auto,
    plain/.style={font=\small, inner sep=2pt}]
    \node[plain] (g)  {$g$};
    \node[plain] (h)  [right of=g]  {$h$};
    \node[plain] (N)  [right of=h]  {$N$};
    \node[plain] (M)  [right of=N]  {$M$};
    \node[plain] (Tg) [right of=M]  {$Tg = M\sin\varphi$};
    \draw[->] (g)  -- (h)  node[midway, above] {\small Step 1};
    \draw[->] (h)  -- (N)  node[midway, above] {\small Step 2};
    \draw[->] (N)  -- (M)  node[midway, above] {\small Riccati};
    \draw[->] (M)  -- (Tg);
\end{tikzpicture}
\caption{Construction of the map $T$ in the swirling case.}
%\label{fig:Tmap}
\end{figure}

\noindent
Then $Tg(0)=Tg(\frac\pi2)=0$ and $\|Tg\|_{C^{1}[0,\frac{\pi}{2}]}\le 3$; hence $T$ maps $\mathfrak{Y}_{3}$ into itself.


For $g_1,g_2\in\mathfrak{Y}_3$, let $h_1,h_2$ be the corresponding solutions of \eqref{eqn:swirlh1}. Lemma~\ref{lm:ODEEst} gives
\[
\|h_i\|_{C^1[0,\frac{\pi}{2}]}\le C(3)\|Z\|_{C^0[0,\frac{\pi}{2}]},\qquad
\|h_1-h_2\|_{C^1[0,\frac{\pi}{2}]}\le C(3)^2\|Z\|_{C^0[0,\frac{\pi}{2}]}\|g_1-g_2\|_{C^0[0,\frac{\pi}{2}]}.
\]
Denote by $N_i, G_i,\mcC_i$ the quantities corresponding to $h_i$.
Repeating the estimates of Step~2, we obtain
\begin{align*}
\|N_1-N_2\|_{C^0[0,1]} 
&\le \frac12 \|(G_1+2\mcC_1)-(G_2+2\mcC_2)\|_{C^0[0,\frac{\pi}{2}]}\\
&\le \frac12 \sup_{0\le\tau,\varphi\le\pi/2} |(G_1-G_2)(\varphi)-(G_1-G_2)(\tau)|\\
&\le \frac{\pi}{4} \|(G_1-G_2)'\|_{C^0[0,\frac{\pi}{2}]},
\end{align*}
where we used \eqref{eqn:NandJprime} and \eqref{eqn:Gand2C}.
The derivative of the difference is
\begin{align*}
(G_1-G_2)' &= 2(h_1h_1'-h_2h_2') + 2(h_1^2-h_2^2)\cot\varphi\\
&=2(h_1-h_2)h_1'+ 2 h_2(h_1-h_2)' + 2(h_1-h_2) (h_1+h_2)\cot\varphi,
\end{align*}
which gives
\[
\|(G_1-G_2)'\|_{C^0[0,\frac{\pi}{2}]}\le 2\pi(\|h_1\|_{C^1[0,\frac{\pi}{2}]}+\|h_2\|_{C^1[0,\frac{\pi}{2}]})\|h_1-h_2\|_{C^1[0,\frac{\pi}{2}]}.
\]
Here we used $|h_1-h_2| \le \frac{\pi}{2} \|(h_1-h_2)'\|_{C^0[0,\frac{\pi}{2}]}$, $|h_2| \le \frac{\pi}{2} \|h_2'\|_{C^0[0,\frac{\pi}{2}]}$, and $|h_i\cot\varphi|\le \|h_i'\|_{C^0[0,\frac{\pi}{2}]}$ for $i=1,2$.
Substituting the $h$-estimates from above,
\[
\|(G_1-G_2)'\|_{C^0[0,\frac{\pi}{2}]}\le 4 C(3)^3 \pi \|Z\|_{C^0[0,\frac{\pi}{2}]}^2\|g_1-g_2\|_{C^0[0,\frac{\pi}{2}]},
\]
and therefore
\[
\|N_1-N_2\|_{C^0[0,1]}
\le C(3)^3  \pi^2 \|Z\|_{C^0[0,\frac{\pi}{2}]}^2\|g_1-g_2\|_{C^0[0,\frac{\pi}{2}]}.
\]

The smallness condition already imposed gives $\|N_{i}\|_{C^{0}}\le\frac14$; Lemma~\ref{lm:Riccati}(ii) then implies
\[
\|M_{1}-M_{2}\|_{C^{1}[0,1]}\le 4\|N_{1}-N_{2}\|_{C^{0}[0,1]}
\le 4 C(3)^{3} \pi^{2} \|Z\|_{C^{0}[0,\frac{\pi}{2}]}^{2}\|g_{1}-g_{2}\|_{C^{0}[0,\frac{\pi}{2}]}.
\]
Since $Tg(\varphi)=M(\cos\varphi)\sin\varphi$, it follows that
\begin{align}\label{eqn:swirlContraction}
\|Tg_1-Tg_2\|_{C^1[0,\frac{\pi}{2}]}\le 12  C(3)^3 \pi^2 \|Z\|_{C^0[0,\frac{\pi}{2}]}^2\|g_1-g_2\|_{C^0[0,\frac{\pi}{2}]}.
\end{align}
By further shrinking $\|Z\|_{C^{0}[0,\frac{\pi}{2}]}$ we can enforce $12\pi^{2} C(3)^{3}\|Z\|_{C^{0}[0,\frac{\pi}{2}]}^{2}<1$,
which makes $T$ a contraction on $\mathfrak{Y}_{3}$. Thus $T$ possesses a unique fixed point $g\in\mathfrak{Y}_{3}$.  Let $h$ be the unique solution of \eqref{eqn:swirlh1} associated with this $g$.  
Define $f = -(g' + g\cot\varphi)$ and recover $P$ from \eqref{eqn:eqnforP}; then $f,g,h\in C^{0}$ and $P\in C^{0}$.  
Exactly as in the swirl-free case, the velocity field $\Bu = \rho^{-1}(f\Be_\rho + g\Be_\varphi + h\Be_\theta)$ and the pressure $p = \rho^{-2}P$ satisfy the Navier--Stokes equations \eqref{eqn:NSEhfD} in the distributional sense.  
Standard Stokes regularity (see \cite[Chapter~IV]{Galdi11}) then yields $\Bu\in C^{2,\alpha}(\R^{3}_{+})$ and $p\in C^{1,\alpha}(\R^{3}_{+})$ for every $\alpha\in(0,1)$.
\end{proof}



\section{Some discussions}\label{sec:Discussions}

In this section we show that the analysis carried out for the half-space extends naturally to a solid cone.  
For $\varphi_0\in(0,\pi)$, let
\begin{equation}\label{eqn:defsolidcone}
\mathcal{K}_{\varphi_0}= \{ (\rho,\theta,\varphi): \rho>0,\; \theta\in(0,2\pi],\; \varphi\in[0,\varphi_0] \}
\end{equation}
denote the solid cone with opening angle $\varphi_0$. 
The stationary Navier--Stokes equations \eqref{eqn:NSEhfD} on $\mathcal{K}_{\varphi_0}$ enjoy the same scaling invariance 
$\Bu(\Bx)\to\lambda\Bu(\lambda\Bx)$, and since $\mathcal{K}_{\pi/2}=\R^3_+$, it is natural to study axisymmetric self-similar solutions on such conical domains.

Inserting
\[
\Bu = \frac{f(\varphi)}{\rho}\Be_\rho + \frac{g(\varphi)}{\rho}\Be_\varphi + \frac{h(\varphi)}{\rho}\Be_\theta,\quad
p = \frac{P(\varphi)}{\rho^2},\quad
\BF = \frac{1}{\rho^3}\bigl(X(\varphi)\Be_\rho + Y(\varphi)\Be_\varphi + Z(\varphi)\Be_\theta\bigr)
\]
into the Navier--Stokes equations yields the same system \eqref{eqn:hDgeneral1} on $[0,\varphi_0]$, but the boundary conditions become
\begin{equation}\label{eqn:coneBC}
f'(0)=g(0)=h(0)=0,\qquad
f(\varphi_0)=g(\varphi_0)=h(\varphi_0)=0.
\end{equation}


\subsection{Reduction of the system.}

The reduction follows that in Sections~\ref{sec:reductNo-swirl} and~\ref{sec:reduct-swirl}, with $\pi/2$ replaced by $\varphi_0$.

If $Z\equiv0$, the same argument as before shows that $h\equiv0$: setting $H=h'+h\cot\varphi$, integrating $H\sin\varphi=(h\sin\varphi)'$ over $[0,\varphi_0]$, and using $h(\varphi_0)=0$ yields $\int_0^{\varphi_0} H\sin\varphi\,\mathrm{d}\varphi=0$, which forces $H\equiv0$ and hence $h\equiv0$.  Thus a swirl-free force produces a swirl-free solution; below we treat the general case $Z\not\equiv0$.

From \eqref{eqn:hDgeneral1}$_{1,2,4}$ we obtain, on $[0,\varphi_0]$,
\[
-f'\sin\varphi + fg\sin\varphi + 2g\sin\varphi
= \int_0^{\varphi} \bigl( G(\tau) + 2\mcC_{\varphi_0} \bigr) \sin\tau \,\rmd\tau,
\]
with $G(\varphi) = X(\varphi) + h^2(\varphi) + 2\int_0^{\varphi} \bigl( Y(\tau) + h^2(\tau)\cot\tau \bigr) \,\rmd\tau$ and a constant $\mcC_{\varphi_0}$.
Letting $L(t)=g(\varphi)\sin\varphi$ on $[t_0,1]$ and $J(t)=\int_0^\varphi (G+2\mcC_{\varphi_0})\sin\tau\,\rmd\tau$, where $t=\cos\varphi$ and $t_0=\cos\varphi_0$, the equation becomes
\[
(1-t^2)\frac{\mathrm{d}^2 L}{\mathrm{d}t^2} + L\frac{\mathrm{d}L}{\mathrm{d}t} + 2L = J \quad\text{on } [t_0,1].
\]
The boundary conditions \eqref{eqn:coneBC} read $L(t_0)=L(1)=\frac{\mathrm{d}L}{\mathrm{d}t}(t_0)=0$.  Integrating from $t$ to $1$ gives
\[
(1-t^2)\frac{\mathrm{d}L}{\mathrm{d}t} + \frac12 L^2 + 2tL = -\int_t^1 J(s)\,\rmd s,
\]
and evaluating at $t=t_0$ forces $\int_{t_0}^1 J(s)\,\rmd s = 0$, which determines
\begin{equation}\label{eq:coneC}
\mcC_{\varphi_0} = -\frac{1}{(1-\cos\varphi_0)^2}\int_0^{\varphi_0} G(\tau)\sin\tau\,(\cos\tau-\cos\varphi_0)\,\rmd\tau .
\end{equation}
Defining $M(t)=L(t)/(1-t^2)$ yields the Riccati equation
\begin{equation}\label{eq:coneRiccati}
\frac{\mathrm{d}M}{\mathrm{d}t} + \frac12 M^{2} = N,\quad t\in[t_0,1],\qquad M(t_0)=0,
\end{equation}
where $N(t)=-\frac{1}{(1-t^2)^2}\int_t^1 J(s)\,\rmd s$ is continuous on $[t_0,1]$ since $\lim\limits_{t\to1^-}N(t)=-\frac18(X(0)+2\mcC_{\varphi_0})$.
The swirl component $h$ satisfies
\begin{equation}\label{eqn:coneSwirlh}
(h'+h\cot\varphi)' = g\,(h'+h\cot\varphi) - Z,\quad \varphi\in[0,\varphi_0],\qquad h(0)=h(\varphi_0)=0,
\end{equation}
with $g(\varphi)=M(\cos\varphi)\sin\varphi$, thereby coupling the two equations.

If $\BF\equiv0$, then $Z\equiv0$ and the argument above gives $h\equiv0$; hence $G\equiv0$ and $\mathcal{C}_{\varphi_0}=0$ by \eqref{eq:coneC}.  Consequently $N\equiv0$ in \eqref{eq:coneRiccati}, and the Riccati equation reduces to $\frac{\mathrm{d}M}{\mathrm{d}t}+\frac12 M^{2}=0$ with $M(t_{0})=0$, whose unique solution is $M\equiv0$.  Thus $f=g=0$ and $P\equiv0$, and we obtain the following rigidity result.

\begin{corollary}
The only axisymmetric self-similar solution of \eqref{eqn:NSEhfD} on $\mathcal{K}_{\varphi_0}$ with $\BF\equiv0$ is $\Bu\equiv0$.
\end{corollary}


\subsection{Existence and uniqueness for small forces.}
We now give a concise proof of the existence theorem on a solid cone, following the same strategy as in Theorem~\ref{thm:A}.
Set $\ell = 1-\cos\varphi_0$ and $t_0 = \cos\varphi_0$.

Similar to the proof of Theorem~\ref{thm:A}, we need an estimate of $N$ in terms of $G'$. In the cone setting, let $t = \cos\varphi$, so that $\varphi\in[0,\varphi_0]$ corresponds to $t\in[t_0,1]$.  Define
\[
N(t) = -\frac{1}{(1-t^{2})^{2}}\int_{t}^{1} J(s)\,\mathrm{d}s,\qquad
J(t) = \int_{0}^{\varphi} \bigl(G(\tau)+2\mathcal{C}_{\varphi_0}\bigr)\sin\tau\,\mathrm{d}\tau,
\]
with $\mathcal{C}_{\varphi_0}$ given by \eqref{eq:coneC}.
Introduce the weight $w(\tau) = \frac{2}{\ell^{2}}\sin\tau(\cos\tau-\cos\varphi_0) \ge 0$, which satisfies $\int_{0}^{\varphi_0}w(\tau)\,\mathrm{d}\tau = 1$.
With this weight we obtain the representation
\begin{equation}\label{eqn:cone-Gand2C}
G(\varphi)+2\mathcal{C}_{\varphi_0} = \int_{0}^{\varphi_0} \bigl(G(\varphi)-G(\tau)\bigr) w(\tau)\,\mathrm{d}\tau .
\end{equation}
Applying \eqref{eqn:NandJprime} on $[t_{0},1]$ and using \eqref{eqn:cone-Gand2C} yields
\[
\|N\|_{C^0[t_0,1]} \le \frac{1}{2(1+t_{0})^{2}}\,\|G+2\mathcal{C}_{\varphi_0}\|_{C^{0}[0,\varphi_0]}
\le \frac{\varphi_{0}}{2(1+t_{0})^{2}}\|G'\|_{C^{0}[0,\varphi_0]} .
\]

\noindent\textbf{Case 1: Swirl-free external force.}
Here $\BF=\rho^{-3}(X\Be_\rho+Y\Be_\varphi)$ and $G=X+2\int_0^{\varphi} Y$. Equation \eqref{eqn:sphericalcurl} gives
\[
\|(\curl\,\BF)^{\tan}\|_{C^0(\mathbb{S}^{2}\cap\mathcal{K}_{\varphi_0})} = \|X'+2Y\|_{C^0[0,\varphi_0]} = \|G'\|_{C^0[0,\varphi_0]}.
\]
By Lemma~\ref{lm:Riccati}(i), a unique solution $M\in C^{1}[t_0,1]$ exists provided $\|N\|_{C^{0}}<\frac{1}{2\ell^{2}}$, which is guaranteed by
\begin{equation}\label{eqn:cone-nos-exist}
\|(\curl\,\BF)^{\tan}\|_{C^0(\mathbb{S}^{2}\cap\mathcal{K}_{\varphi_0})} < \frac{(1+\cos\varphi_0)^{2}}{\varphi_0\,(1-\cos\vp_0)^{2}} =: \varepsilon_{\mathrm{sff}}(\varphi_0).
\end{equation}
When this holds, $g(\varphi) = M(\cos\varphi)\sin\varphi \in C^1[0,\varphi_0]$ and $f = -(g' + g\cot\varphi) \in C^0[0,\varphi_0]$, and the Stokes regularity theory gives
$\Bu\in C^{2,\alpha}(\mathcal{K}_{\varphi_0})$, $p\in C^{1,\alpha}(\mathcal{K}_{\varphi_0})$ for every $\alpha \in (0,1)$.

\noindent\textbf{Case 2: Swirling external force.}
Now $\BF=\rho^{-3}(X\Be_\rho+Y\Be_\varphi+Z\Be_\theta)$. Equation \eqref{eqn:sphericalcurl} shows
\[
\|(\curl\,\BF)^{\tan}\|_{C^0(\mathbb{S}^{2}\cap\mathcal{K}_{\varphi_0})} = \bigl\|2Z\Be_{\varphi} - (X'+2Y)\Be_{\theta}\bigr\|_{C^0(\mathbb{S}^2 \cap \mathcal{K}_{\varphi_0})},\quad \text{comparable to }2\|Z\|_{C^0[0,\varphi_0]} + \|X'+2Y\|_{C^0[0,\varphi_0]}.
\]
The proof follows the same contraction argument used for the swirling external force in Section~\ref{sec:ExistSwirl}.
Let
$\mathfrak{Y}_3 = \{ g\in C^1[0,\varphi_0] : g(0)=g(\varphi_0)=0,\ \|g\|_{C^1[0,\varphi_0]}\le 3 \}$.
For $g\in\mathfrak{Y}_3$, let $h$ be the unique solution of \eqref{eqn:coneSwirlh} supplied by Lemma~\ref{lm:ODEEst}. Then
\[
\|h\|_{C^1[0,\varphi_0]} \le C_{\mathrm{h}}(\varphi_0)\,\|Z\|_{C^0[0,\varphi_0]},
\]
where $C_{\mathrm{h}}$ remains bounded as $\varphi_0\to0$.
Using $G$ as in \eqref{eqn:coneG} and replacing $\pi/2$ by $\varphi_0$ in the Step~2 estimates of Section~\ref{sec:ExistSwirl} yields
\[
\|G'\|_{C^0[0,\varphi_0]} \le \|X'+2Y\|_{C^0[0,\varphi_0]} + \varphi_0\,C_{\rm h}(\varphi_0)^2\,\|Z\|_{C^0[0,\varphi_0]}^{2}.
\]
Consequently
\begin{equation}\label{eq:coneNest}
\|N\|_{C^0[t_0,1]} \le \frac{\varphi_0}{2(1+\cos\varphi_0)^{2}}\,\|X'+2Y\|_{C^0[0,\varphi_0]}
+ \frac{\varphi_0^{2}}{2(1+\cos\varphi_0)^{2}}\,C_{\rm h}(\varphi_0)^2\,\|Z\|_{C^0[0,\varphi_0]}^{2}.
\end{equation}
%where $C_1$ is bounded as $\varphi_0\to0$.

\emph{Self-mapping.}  
If $\|N\|_{C^0[t_0,1]}<\frac{1}{2\ell^{2}}$, Lemma~\ref{lm:Riccati}(i) gives a unique solution $M\in C^1[t_0,1]$ with $\|M\|_{C^1[t_0,1]}\le 2(\ell+1)\|N\|_{C^0[t_0,1]}$.  
Define $Tg(\varphi):=M(\cos\varphi)\sin\varphi$; a direct computation shows
\[
\|Tg\|_{C^1[0,\vp_0]}\le 3\|M\|_{C^1[t_0,1]}\le 6(\ell+1)\|N\|_{C^0[t_0,1]}.
\]
Thus $Tg\in\mathfrak{Y}_3$ as soon as $\|N\|_{C^0[t_0,1]}\le \frac{1}{2(\ell+1)}$.  
To prepare for the subsequent contraction argument we enforce the slightly stronger bound $\|N\|_{C^0[t_0,1]}\le\frac{1}{4\ell^{2}}$.  
By \eqref{eq:coneNest}, taking $\|X'+2Y\|_{C^0[0,\vp_0]}$ and $\|Z\|_{C^0[0,\vp_0]}$ sufficiently small guarantees
\[
\|N\|_{C^0[t_0,1]}\le \min\Bigl\{\frac{1}{2(\ell+1)},\,\frac{1}{4\ell^{2}}\Bigr\},
\]
and therefore $T$ maps $\mathfrak{Y}_3$ into itself.

\emph{Contraction.}  Repeating the difference estimates of Step~3 in Section~\ref{sec:ExistSwirl} on $[0,\varphi_0]$ and using Lemma~\ref{lm:ODEEst}(ii), we obtain
\[
\|Tg_1-Tg_2\|_{C^1[0,\vp_0]} \le K(\vp_0)\varphi_0^{2}\,\|Z\|_{C^0[0,\vp_0]}^{2}\,\|g_1-g_2\|_{C^0[0,\vp_0]}.
\]
Here the constant $K(\vp_0)$ remains bounded as $\varphi_0\to0$. Consequently, $T$ is a contraction on $\mathfrak{Y}_3$ provided $\|Z\|_{C^0[0,\varphi_0]} \le  (\sqrt{2\,K(\varphi_0)} \vp_0)^{-1}$.

Therefore, there exists a threshold $\varepsilon_{\mathrm{sf}}(\varphi_0)>0$ such that whenever
\[
\|X'+2Y\|_{C^0[0,\varphi_0]} + 2\|Z\|_{C^0[0,\varphi_0]} < \varepsilon_{\mathrm{sf}}(\varphi_0),
\]
the map $T$ maps $\mathfrak{Y}_3$ into itself and is a contraction on $\mathfrak{Y}_3$.
Hence $T$ admits a unique fixed point $g\in\mathfrak{Y}_3$, which together with the associated $h$ and $M$ solves the coupled ODE system.
By the Stokes regularity theory we obtain $\Bu\in C^{2,\alpha}(\mathcal{K}_{\varphi_0})$ and $p\in C^{1,\alpha}(\mathcal{K}_{\varphi_0})$ for every $\alpha\in(0,1)$.

From the above discussion we immediately obtain the following existence theorem on a solid cone.

\begin{theorem}[Existence on a solid cone]\label{thm:A-cone}
Let $\varphi_0\in(0,\pi)$ and let $\mathcal{K}_{\varphi_0}$ be the solid cone defined in \eqref{eqn:defsolidcone}.
Let $\BF$ be an axisymmetric vector field of class $C^1$, homogeneous of degree $-3$ on $\mathcal{K}_{\vp_0}$, and let $(\curl\,\BF)^{\tan}$ denote the tangential component of $\curl\,\BF$ on the unit sphere $\mathbb{S}^2$.
Then there exists $\varepsilon(\vp_0)>0$ such that if
\[
\|(\curl\,\BF)^{\tan}\|_{C^0(\mathbb{S}^2\cap\mathcal{K}_{\varphi_0})} \le \varepsilon(\varphi_0),
\]
the system \eqref{eqn:NSEhfD} admits a unique axisymmetric self-similar solution, small in $C^0(\mathcal{K}_{\varphi_0})$, with $\Bu\in C^{2,\alpha}(\mathcal{K}_{\varphi_0})$ and $p\in C^{1,\alpha}(\mathcal{K}_{\varphi_0})$ for every $\alpha\in(0,1)$.

Moreover, if $\BF$ has no swirl, then any axisymmetric self-similar solution is necessarily swirl-free and unique.
\end{theorem}

\begin{remark}
For a swirl-free external force, we have 
$\varepsilon_{\mathrm{sff}}(\varphi_0)=\frac{(1+\cos\varphi_0)^2}{\varphi_0\,(1-\cos\varphi_0)^{2}}\sim 16\,\varphi_0^{-5}$ as $\varphi_0\to0$.
In the swirling case, for small $\varphi_0$ the self-mapping and contraction requirements together force both $\|X'+2Y\|_{C^0[0,\vp_0]}$ and $\|Z\|_{C^0[0,\vp_0]}$ to be at most of order $\varphi_0^{-1}$; consequently $\varepsilon_{\mathrm{sf}}(\varphi_0)\sim c\,\varphi_0^{-1}$.
Thus, in either case, $\varepsilon(\varphi_0)\to\infty$ as $\varphi_0\to0$, so for a sufficiently narrow cone the smallness condition on $\|(\curl\,\BF)^{\tan}\|_{C^0(\mathbb{S}^2\cap\mathcal{K}_{\varphi_0})}$ is weakened,
and a unique solution exists even for a very large tangential curl.
\end{remark}

% \begin{remark}
% \begin{enumerate}[(i)]
% \item If the smallness condition fails on a given cone $\mathcal{K}_{\varphi_1}$, we can simply restrict the same force to a narrower cone $\mathcal{K}_{\varphi_2}$ with $\varphi_2<\varphi_1$. Since $\varepsilon(\varphi_0)\to\infty$ as $\varphi_0\to0$, we can always choose $\varphi_2$ small enough so that the condition is satisfied on $\mathcal{K}_{\varphi_2}$, and Theorem~\ref{thm:A-cone}
% yields a unique self-similar solution there.
% \end{enumerate}
% \end{remark}


\subsection{Blow-up criterion.}
By repeating the arguments of Section~\ref{sec:blowup} line by line, which rely only on the Riccati equation on a general interval $[t_0,1]$, we immediately obtain the following blow-up criterion on a solid cone.

\begin{theorem}[Blow-up criterion on a solid cone]
Let $\varphi_0\in(0,\pi)$ and let $\BF$ be an axisymmetric, swirl-free external force on $\mathcal{K}_{\varphi_0}$, of class $C^{1}$ and homogeneous of degree $-3$.
There exist constants $\underline{\lambda}<0<\overline{\lambda}$ (possibly $\pm\infty$) with the following property: for every $\lambda\in(\underline{\lambda},\overline{\lambda})$ the Navier--Stokes equations \eqref{eqn:NSEhfD} with $\mathbb{R}^{3}_{+}$ replaced by $\mathcal{K}_{\varphi_0}$ and with external force $\lambda\BF$ admit a unique axisymmetric self-similar solution
$\Bu\in C^{2,\alpha}_{\mathrm{loc}}(\overline{\mathcal K_{\varphi_0}}\setminus\{0\})$, $p\in C^{1,\alpha}_{\mathrm{loc}}(\overline{\mathcal K_{\varphi_0}}\setminus\{0\})$ for every $\alpha\in(0,1)$, whereas no such bounded-profile solution exists for $\lambda\notin(\underline{\lambda},\overline{\lambda})$.

In spherical coordinates we write $\BF = \rho^{-3}\bigl( X(\varphi)\Be_{\rho} + Y(\varphi)\Be_{\varphi} \bigr)$.
Then the finiteness of $\underline{\lambda}$ and $\overline{\lambda}$ is determined by the sign of the discriminant
\begin{align*}
\mathcal{N}_{\vp_0}(\varphi)=\int_{0}^{\varphi}
\bigl(G(\tau)+2\mathcal{C}_{\varphi_0}\bigr)\,\sin\tau\,(\cos\varphi-\cos\tau)\,\rmd\tau,
\quad \varphi \in [0,\varphi_0],
\end{align*}
where $G(\varphi)=X(\varphi)+2\int_{0}^{\varphi} Y(\tau)\,\rmd\tau$ and $\mathcal{C}_{\varphi_0}$ is defined by \eqref{eq:coneC}. More precisely:
\begin{itemize}
\item if $\mathcal{N}_{\vp_0}\ge0$ and $\mathcal{N}_{\vp_0}\not\equiv0$, then $\underline{\lambda}\in(-\infty,0)$ and $\overline{\lambda}=+\infty$;

\item if $\mathcal{N}_{\vp_0}\le0$ and $\mathcal{N}_{\vp_0}\not\equiv0$, then $\underline{\lambda}=-\infty$ and $\overline{\lambda}\in(0,+\infty)$;

\item if $\mathcal{N}_{\vp_0}$ changes sign, then $\underline{\lambda}\in(-\infty,0)$ and $\overline{\lambda}\in(0,+\infty)$;
\item if $\mathcal{N}_{\vp_0}\equiv0$, then $\underline{\lambda}=-\infty$ and $\overline{\lambda}=+\infty$.
\end{itemize}
\end{theorem}

\medskip

{\bf Acknowledgments}The research of Wang is partially supported by NSFC grants 12671279, the Natural Science Foundation of Jiangsu Province (Grant No. BK20240147) and Jiangsu Provincial Scientific Research Center of Applied Mathematics(No. BK20233002). The research of Xie is partially supported by  NSFC grants 12571238, 12426203, and 12250710674. 


\medskip
{\bf Statement and declaration.} The authors have no relevant financial or non-financial interests to disclose. No datasets were generated or analyzed during the current study.




\bibliographystyle{alpha}
\bibliography{ref}

\end{document}
